% ============================================================================
%  CAPÍTULO 3 — METODOLOGÍA
% ============================================================================
\chapter{Metodología}
\label{cap:metodologia}

El diseño metodológico de esta investigación sigue el marco CRISP-DM, cuya
justificación conceptual se desarrolló en la Sección~\ref{sec:crisp-dm}. La
elección de este marco se fundamenta en tres características del problema que lo
hacen especialmente adecuado. Los datos GRD son registros administrativos
concebidos originalmente con fines financieros, presentan heterogeneidad técnica
entre años y carecen de un etiquetado funcional previo validado, lo que exige
dedicar una parte sustancial del proceso a la comprensión del dominio, la
depuración y la validación de consistencia antes de cualquier modelamiento. Al
estructurar estas fases de manera explícita y exigir su documentación, CRISP-DM
permite rastrear cómo una decisión adoptada durante la limpieza repercute en los
perfiles finales. A esto se suma que el objetivo  no es optimizar un
predictor sino describir la configuración funcional de la red hospitalaria, de
modo que la evaluación no puede apoyarse únicamente en la clasificación
administrativa del MINSAL, que sirve como referencia de contraste pero no como
criterio de validación directa de la segmentación funcional, al estar basada
en infraestructura y no en casuística clínica. Finalmente, al separar conceptualmente las etapas de
preparación, modelado y evaluación, el marco facilita comparar distintas
configuraciones del pipeline sin mezclar las decisiones propias de cada fase.

Para situar cada sección de este capítulo dentro del marco, las seis etapas de
CRISP-DM se corresponden con el trabajo desarrollado de la siguiente forma. La
comprensión del negocio se aborda en la sección siguiente,
retomando y precisando los objetivos e hipótesis del Capítulo~\ref{cap:introduccion}
desde la perspectiva del diseño analítico. La comprensión de los datos cubre la
exploración de los archivos GRD anuales, la detección de heterogeneidad técnica
entre años y la identificación de problemas de calidad. La preparación de los
datos abarca la estandarización, limpieza y transformación de los registros, la
construcción de variables institucionales, la reducción dimensional de la
casuística y el preprocesamiento previo al modelado. El modelado corresponde a la
ejecución del clustering jerárquico aglomerativo en dos niveles, incluyendo la
elección del algoritmo, la selección del número de grupos y la subdivisión del
núcleo mayoritario. La evaluación comprende la validación estadística de los
grupos, el submuestreo para estabilidad, la selección de variables relevantes, el
contraste con la clasificación administrativa del MINSAL y la detección de
establecimientos atípicos. Por último, el despliegue se materializa, dentro del
alcance de este trabajo, como la presentación de los perfiles funcionales y su
contraste con la clasificación administrativa vigente, en un formato interpretable
y utilizable por gestores sanitarios.

\section{Herramientas de desarrollo}
\label{sec:herramientas}

El trabajo se implementa en lenguaje Python 3.11\footnote{El desarrollo y
ejecución se realizan en Linux Mint 22.3 (base Ubuntu 24.04) sobre equipo personal
con procesador AMD Ryzen AI 7 350, 24 GB RAM y SSD NVMe. No se requiere
infraestructura de servidores adicional.}, elegido porque permite integrar en un
mismo entorno el procesamiento de datos, el modelado estadístico y la generación
de visualizaciones, y porque las librerías de clustering, validación e
interpretación empleadas están disponibles de forma nativa en este lenguaje.
Las versiones utilizadas y sus roles en el pipeline se detallan en la
Tabla~\ref{tab:software}.

\begin{table}[H]
\centering
\caption{Librerías principales del pipeline y sus roles}
\label{tab:software}
\begin{tabular}{lll}
\hline
\textbf{Librería} & \textbf{Versión} & \textbf{Rol en el pipeline} \\
\hline
\texttt{pandas}      & 2.2.3  & Manipulación y transformación de datos \\
\texttt{numpy}       & 1.26.4 & Cálculo numérico \\
\texttt{scikit-learn}& 1.4.2  & Clustering, métricas de validación y reducción de dimensionalidad \\
\texttt{scipy}       & 1.13.1 & Pruebas estadísticas de contraste entre grupos \\
\texttt{statsmodels} & 0.14.2 & Modelos estadísticos complementarios \\
\texttt{matplotlib}  & 3.8.4  & Visualizaciones \\
\texttt{seaborn}     & 0.13.2 & Visualizaciones estadísticas \\
\texttt{pyarrow}     & 16.1.0 & Almacenamiento intermedio en formato Parquet \\
\hline
\end{tabular}
\end{table}

El pipeline se organiza en una estructura modular con tres componentes que
garantizan reproducibilidad. El paquete \texttt{src/} encapsula la lógica
reutilizable en tres subpaquetes: \texttt{src/etl/} (mapeadores CIE-10 y
CIE-9-MC, construcción de indicadores, diversidad y variables extendidas),
\texttt{src/modeling/} (ensamblado de la matriz institucional, reducción
dimensional, métodos de agrupamiento y comparación de particiones) y
\texttt{src/utils/} (entrada/salida con validación de esquema y orden
determinista). Los \textit{scripts} de \texttt{scripts/} ejecutan el pipeline de
extremo a extremo de forma no interactiva: \texttt{build\_extended\_grd.py}
consolida los archivos GRD anuales, \texttt{build\_dataset.py} construye la
matriz institucional y el barrido de métricas Ward,
\texttt{build\_hierarchical\_clustering.py} fija la partición reportada y la
familia \texttt{analyze\_*.py} y \texttt{compare\_*.py} produce los contrastes
estadísticos y los análisis de sensibilidad. Los \textit{notebooks} Jupyter
(\texttt{01\_etl.ipynb}, \texttt{02\_eda.ipynb}, \texttt{03\_clustering.ipynb},
\texttt{04\_xai\_outliers.ipynb}, \texttt{05\_clustering\_alternativos.ipynb})
conservan el registro exploratorio del análisis. Los archivos
\texttt{requirements.txt} y \texttt{requirements-dev.txt} declaran las versiones
exactas de todas las librerías empleadas, y las semillas aleatorias están
fijadas en el código (\texttt{RANDOM\_STATE = 42}), de modo que cualquier
ejecución posterior reproduzca los mismos resultados. La verificación del
comportamiento de los módulos se automatiza con \texttt{pytest} y
\texttt{hypothesis} en el directorio \texttt{tests/}.

Los datos fuente corresponden a los seis archivos GRD públicos del MINSAL para el
período 2019--2024, disponibles en la plataforma Tablero GRD del Ministerio de
Salud~\citep{TableauGRD}. El código completo del proyecto, incluyendo notebooks,
scripts, pruebas y documentación, está disponible en
\url{https://github.com/reii23/chile-hospital-grd-profiling} bajo licencia MIT,
que permite su reutilización y continuidad conservando la atribución de autoría.
La estructura del repositorio y las condiciones de reproducción se detallan en el
Anexo~\ref{anexo:repositorio}.

\section{Comprensión del negocio}
\label{sec:comprension-negocio}

La etapa de comprensión del negocio en CRISP-DM tiene por propósito traducir
el problema institucional en condiciones de éxito verificables y en restricciones
operativas que orienten todas las decisiones metodológicas posteriores. Dado que
los objetivos, la hipótesis y la descripción del problema se desarrollaron en el
Capítulo~\ref{cap:introduccion}, esta sección se concentra en los aspectos que
son propios del diseño analítico y que no quedaron explicitados allí.

En cuanto al criterio de éxito, este opera en dos dimensiones que se
complementan entre sí. La primera es de naturaleza técnica y exige que los grupos
resultantes presenten cohesión interna y separación mutua estadísticamente
sustentables, evaluadas mediante el coeficiente de Silhouette, el índice de
Calinski-Harabasz y el índice de Davies-Bouldin, respaldadas por pruebas de
estabilidad a través de remuestreo. La segunda dimensión es de dominio y exige
que cada perfil pueda caracterizarse con atributos de casuística y producción que
tengan sentido desde la gestión sanitaria, más allá de constituir una partición
matemáticamente óptima. Ambas dimensiones se complementan con una condición
adicional: los perfiles obtenidos deben exhibir mayor homogeneidad interna que la
clasificación administrativa del MINSAL en las variables de producción clínica,
lo que se verifica mediante el contraste estadístico entre ambas segmentaciones.

En lo que respecta a las restricciones operativas, el proyecto se enmarca en
tres definiciones que delimitan su alcance. La unidad de análisis es el hospital
y no el paciente individual, por lo que todos los indicadores se construyen de
forma agregada a nivel institucional a partir de los egresos. El análisis opera
exclusivamente sobre los datos disponibles en la base GRD, excluyendo información
de costos individuales y dotación de camas, de manera que el perfilamiento se
fundamente en la actividad clínica registrada y no en la capacidad instalada.
Por último, la reproducibilidad del proceso es un requisito transversal a todas
las etapas, lo que implica que cada decisión metodológica debe quedar documentada
y ser verificable, permitiendo replicar y actualizar el pipeline cuando se
incorporen nuevos años de datos al sistema.

Esta definición de la unidad de análisis tiene una consecuencia directa sobre el
tamaño muestral relevante para toda la inferencia estadística del trabajo. Si
bien la base procesada contiene 5{,}808{,}515 egresos individuales, estos se
agregan en indicadores institucionales, de modo que el clustering, las pruebas
de significancia, la detección de atípicos y el contraste con la clasificación
MINSAL operan sobre $n = 65$ hospitales, no sobre el volumen de egresos. El
tamaño de la base de egresos permite construir indicadores institucionales
estables (Sección~\ref{sec:elegibilidad}), pero no debe interpretarse como una
fuente de alta potencia estadística para el análisis multivariado, que se
realiza a nivel de establecimiento. Esta distinción es aún más relevante en el
Nivel~2 (Sección~\ref{sec:nivel2-metodo}), donde el tamaño muestral efectivo se
reduce a $n = 54$ hospitales y algunos subgrupos resultantes contienen apenas
una o dos observaciones, lo que exige cautela en la interpretación de cualquier
prueba estadística aplicada a esos subgrupos.

\section{Comprensión de los datos}
\label{sec:datos}

La fuente primaria de este estudio son los archivos públicos de egresos
hospitalarios codificados con el sistema GRD, provistos por el Ministerio de
Salud de Chile para el período 2019--2024. Los datos se encuentran disponibles
en formato \texttt{.txt} y se distribuyen en seis archivos anuales, uno por año.
La Tabla~\ref{tab:volumen-registros} presenta el volumen de registros originales
de cada archivo antes de cualquier proceso de limpieza o filtrado.

\begin{table}[htbp]
\centering
\caption{Volumen de registros originales por archivo anual GRD.}
\label{tab:volumen-registros}
\begin{tabular}{lll}
\toprule
\textbf{Año} & \textbf{Archivo} & \textbf{Registros} \\
\midrule
2019 & \texttt{GRD\_PUBLICO\_2019.txt} & 1{,}151{,}475 \\
2020 & \texttt{GRD\_PUBLICO\_2020.txt} &   781{,}912 \\
2021 & \texttt{GRD\_PUBLICO\_2021.txt} &   816{,}909 \\
2022 & \texttt{GRD\_PUBLICO\_EXTERNO\_2022.txt} &   932{,}840 \\
2023 & \texttt{GRD\_PUBLICO\_2023.txt} & 1{,}039{,}587 \\
2024 & \texttt{GRD\_PUBLICO\_2024.txt} & 1{,}085{,}813 \\
\midrule
\textbf{Total} & & \textbf{5{,}808{,}536} \\
\bottomrule
\end{tabular}
\end{table}

Cada registro corresponde a un episodio de egreso hospitalario e incorpora
información organizada en cuatro dimensiones. La primera comprende datos del
paciente, entre los que se encuentran el sexo, la fecha de nacimiento y la
comuna de residencia. La segunda abarca información hospitalaria, que incluye
las fechas de ingreso y alta, el tipo de actividad, el tipo de ingreso y egreso,
el servicio clínico y el establecimiento tratante. La tercera corresponde a la
clasificación GRD propiamente tal, con el código del grupo asignado, su peso
relativo, el índice de severidad y el índice de mortalidad del episodio. La
cuarta dimensión recoge la codificación clínica estandarizada, con hasta 35
diagnósticos en CIE-10, que permiten identificar la causa principal de
hospitalización y las comorbilidades, y hasta 30 procedimientos en CIE-9-MC,
que registran las intervenciones quirúrgicas y exámenes de mayor complejidad
realizados durante la estadía. Adicionalmente, cada registro puede contener
información complementaria sobre antecedentes neonatales, traslados internos y
uso de pabellón quirúrgico.

La Tabla~\ref{tab:variables-grd} detalla las principales variables del registro,
su descripción y el tipo de dato con que se almacenan en los archivos originales.
Las columnas de traslados (\texttt{FECHATRASLADO1}--\texttt{SERVICIOTRASLADO9})
y los campos de diagnósticos y procedimientos se omiten de la tabla por ser
repetitivos, ya que cada uno sigue la misma estructura que su primer elemento.

Los registros se encuentran completamente anonimizados mediante el campo
\texttt{CIP\_ENCRIPTADO}, cuya reversión no es técnicamente posible, por lo que
su uso es compatible con la Ley N°~19.628 de protección de datos personales sin
requerir consentimiento informado ni evaluación por comité de ética.

\subsection{Heterogeneidad técnica de las fuentes}
\label{sec:heterogeneidad}

Un desafío que emerge al explorar los archivos es la inconsistencia técnica entre
años, que obliga a un tratamiento diferenciado antes de cualquier proceso de
unificación. La Tabla~\ref{tab:heterogeneidad} resume las principales
discrepancias detectadas.

\begin{table}[htbp]
\centering
\caption{Heterogeneidad técnica de los archivos GRD por año.}
\label{tab:heterogeneidad}
\begin{tabular}{lll}
\toprule
\textbf{Año} & \textbf{Codificación} & \textbf{Identificador de paciente} \\
\midrule
2019 & UTF-8        & \texttt{CIP\_ENCRIPTADO} \\
2020 & UTF-8        & \texttt{CIP\_ENCRIPTADO} \\
2021 & UTF-8-SIG    & \texttt{CIP\_ENCRIPTADO} \\
2022 & UTF-16       & \texttt{CIP\_ENCRIPTADO} \\
2023 & UTF-16       & \texttt{CIP\_ENCRIPTADO} \\
2024 & ISO-8859-1   & \texttt{ID\_BENEFICIARIO} \\
\bottomrule
\end{tabular}
\end{table}

Esta variación de codificación no es menor. Leer un archivo con una codificación
distinta a la original corrompe los caracteres acentuados presentes en nombres
de establecimientos y diagnósticos, impidiendo la unificación posterior. Este
problema fue detectado durante la exploración inicial y condicionó directamente
el diseño del proceso de estandarización descrito en la etapa siguiente.

El cambio del identificador de paciente de \texttt{CIP\_ENCRIPTADO}
(2019--2023) a \texttt{ID\_BENEFICIARIO} (2024) se documenta en la
Tabla~\ref{tab:heterogeneidad} por representar una discrepancia real entre
los archivos fuente, pero no tiene consecuencias para este trabajo: la
unidad de análisis es el hospital (Sección~\ref{sec:comprension-negocio}) y
ninguna variable de la matriz institucional se construye a partir de la
identidad de un paciente individual ni de su seguimiento entre años. En
consecuencia, este campo se excluye deliberadamente de las columnas que el
pipeline de extracción lee de los archivos originales, y no se requiere
demostrar que ambos identificadores sean equivalentes ni comparables entre
años para ningún indicador reportado en este trabajo.

La Figura~\ref{fig:pipeline-datos} resume el flujo completo desde los registros
originales hasta la matriz institucional preprocesada que sirve de entrada al
modelado.

\begin{figure}[H]
\centering
\begin{tikzpicture}[
    node distance = 7mm and 0mm,
    box/.style = {
        rectangle, rounded corners=3pt,
        draw=black!70, fill=white,
        text width=8.5cm, align=center,
        minimum height=1.0cm, font=\small
    },
    filter/.style = {
        rectangle, rounded corners=3pt,
        draw=black!40, fill=black!5,
        text width=8.5cm, align=center,
        minimum height=0.8cm, font=\footnotesize\itshape
    },
    arrow/.style = {-{Stealth[length=5pt]}, thick, black!60}
]

% Nodo 1
\node[box, fill=blue!8] (raw)
    {\textbf{Registros originales}\\
     6 archivos GRD 2019--2024 $\;\cdot\;$ \textbf{5{,}808{,}536 registros}};

% Filtro 1
\node[filter, below=of raw] (f1)
    {Estandarización de codificación (UTF-8/UTF-16/ISO-8859-1),
     homologación de columnas y eliminación de 21 registros inválidos};

% Nodo 2
\node[box, below=of f1, fill=blue!8] (clean)
    {\textbf{Registros depurados}\\
     \textbf{5{,}808{,}515 registros} consolidados en un único archivo \texttt{parquet}};

% Filtro 2
\node[filter, below=of clean] (f2)
    {Criterio de elegibilidad: presencia en $\geq 3$ de los 6 años del período\\
     (excluye 7 hospitales incorporados recién en 2023--2024)};

% Nodo 3
\node[box, below=of f2, fill=blue!8] (hosp)
    {\textbf{Hospitales elegibles}\\
     \textbf{65 establecimientos} (de 72 con datos en el período)};

% Filtro 3
\node[filter, below=of hosp] (f3)
    {Construcción de variables institucionales (4 grupos) y
     reducción PCA de la casuística CIE-10 / CIE-9-MC / Top-20 GRD
     (8 componentes)};

% Nodo 4
\node[box, below=of f3, fill=blue!8] (matrix)
    {\textbf{Matriz institucional (\textit{Hospital Matrix Integrada})}\\
     \textbf{65 hospitales $\times$ 35 variables}};

% Filtro 4
\node[filter, below=of matrix] (f4)
    {Preprocesamiento: transformación $\log(1+x)$ de las variables asimétricas
     y escalado robusto (mediana e IQR)};

% Nodo 5
\node[box, below=of f4, fill=green!10] (scaled)
    {\textbf{Matriz preprocesada}\\
     \textbf{65 hospitales $\times$ 35 variables} — entrada al modelado};

% Flechas
\draw[arrow] (raw)   -- (f1);
\draw[arrow] (f1)    -- (clean);
\draw[arrow] (clean) -- (f2);
\draw[arrow] (f2)    -- (hosp);
\draw[arrow] (hosp)  -- (f3);
\draw[arrow] (f3)    -- (matrix);
\draw[arrow] (matrix) -- (f4);
\draw[arrow] (f4)    -- (scaled);

\end{tikzpicture}
\caption{Flujo de datos desde los registros originales hasta la matriz institucional preprocesada.}
\label{fig:pipeline-datos}
\end{figure}

\subsection{Limitaciones de la fuente}
\label{sec:limitaciones-fuente}

Por tratarse de registros administrativos concebidos con fines de financiamiento
y no de investigación, la base GRD presenta tres limitaciones que deben
considerarse al interpretar los resultados del perfilamiento.

La primera tiene que ver con la exhaustividad diagnóstica. Los hospitales con
mayor dotación de codificadores o con sistemas de apoyo a la codificación tienden
a registrar más diagnósticos secundarios por egreso que aquellos con menor
capacidad técnica, lo que puede inflar artificialmente indicadores como el
promedio de comorbilidades en ciertos establecimientos. Aunque al agregar los
registros a nivel institucional este efecto se distribuye entre todos los egresos
del hospital y pierde parte de su influencia, no desaparece por completo. En
consecuencia, una diferencia entre hospitales en las variables de comorbilidad no
admite una lectura directa como diferencia en la carga clínica de los pacientes
atendidos, porque puede originarse en la práctica de registro de cada centro.

Esta limitación merece una precisión cuantitativa, porque la decisión de leer
la totalidad de los campos de codificación (Sección~\ref{sec:variables}) la
vuelve más visible en lugar de atenuarla. Los 65 hospitales elegibles utilizan
campos más allá del quinto diagnóstico, con una profundidad que va de $0{,}73$ a
$3{,}80$ diagnósticos adicionales por egreso, de modo que la variación en
exhaustividad no afecta a un subconjunto marginal sino a toda la red. Al
descomponer esa variación, el Servicio de Salud al que pertenece el
establecimiento explica un $51\,\%$ de la varianza en profundidad de
codificación, proporción que sugiere un componente institucional o
administrativo y no puramente clínico. Dos elementos acotan esa lectura. Por
una parte, la dispersión dentro de un mismo Servicio de Salud sigue siendo
amplia --en Araucanía Sur la profundidad va de $1{,}28$ a $3{,}80$--, por lo que
el fenómeno no se reduce a un patrón regional homogéneo. Por otra, el aumento
del promedio de comorbilidades al incorporar los campos completos correlaciona
positivamente con la severidad media ($r = 0{,}64$) y con la mortalidad media
($r = 0{,}58$) de cada establecimiento, es decir, se concentra donde la carga
clínica registrada es mayor, lo que es coherente con un componente clínico
genuino. Conviene subrayar que truncar los campos no controlaba este sesgo sino
que lo encubría: bajo un techo de cinco diagnósticos, un hospital que codifica
veinte y otro que codifica cinco resultaban indistinguibles. La cobertura
completa hace la heterogeneidad de registro observable y, por lo mismo,
obligatoria de declarar al interpretar cualquier comparación de comorbilidad
entre establecimientos.

La segunda limitación se relaciona con la estabilidad longitudinal del sistema
GRD y de la producción hospitalaria a lo largo del período. Esta limitación opera
en tres dimensiones.

La primera es la pandemia de COVID-19. Los años 2020 y 2021 presentan caídas
bruscas en el volumen de registros, con 781{,}912 y 816{,}909 egresos
respectivamente, frente a 1{,}151{,}475 en 2019, lo que refleja la contracción de
la actividad electiva durante la crisis sanitaria. Esta disminución no fue uniforme entre
hospitales: los centros de mayor complejidad reconvirtieron actividad hacia la
atención de pacientes críticos, alterando temporalmente sus perfiles de
casuística, índices de severidad, estancias medias y tasas de mortalidad
registrada. Al calcular promedios institucionales sobre los seis años, los
valores de 2020--2021 quedan incluidos con el mismo peso que los años de
actividad normal, lo que puede atenuar diferencias funcionales reales entre
establecimientos. Este efecto no puede eliminarse del pipeline de análisis, y
constituye una razón adicional para interpretar los perfiles como una
caracterización del período en su conjunto y no como representación de la
actividad habitual pre o post pandemia.

La segunda dimensión es la transición del manual que rige la captura y el
procesamiento de los egresos. Los años 2019 y 2020 operaron bajo la versión 4.0
del Manual de Orientación para la Captura y Procesamiento de los Egresos
Hospitalarios, mientras que los años 2022 a 2024 lo hicieron bajo la versión 5.0,
aprobada mediante Resolución Exenta N°~934 del 14 de diciembre de 2021 e
instruida a partir de 2022, correspondiendo el año 2021 a un período de
transición. El alcance real de esta discontinuidad se acotó inspeccionando los
propios archivos. Los seis años reportan el mismo agrupador (IR-GRD 29301) y,
más importante, la relación entre código GRD y peso relativo es constante en
todo el período: entre los 1.064 códigos GRD distintos presentes en la base
depurada, ninguno registra más de un valor de peso relativo ni más de un valor
de índice de severidad a lo largo de los seis años. Es decir, los archivos
públicos emplean una única tabla de pesos e índices, y la transición de versión
no introdujo un reescalamiento de estos atributos. Lo que no puede verificarse
desde los archivos es la asignación misma de los episodios a grupos: un mismo
cuadro clínico podría quedar clasificado en códigos distintos según la versión
vigente, y el conjunto de códigos GRD observados varía moderadamente entre años
(de 1.053 en 2019 a 989 en 2024). Ese margen de incertidumbre, y no una
diferencia de pesos, es el ruido que la transición introduce en los indicadores
de complejidad que cruzan el cambio.

La tercera dimensión es normativa: a partir del 30 de enero de 2023, la Ley de
Presupuestos N°~21.516 exigió registrar bajo estándares GRD el 100\,\% de la
actividad de atención abierta, incluyendo la hospitalización diurna, lo que puede
alterar la composición de los registros de 2023 y 2024 respecto a años anteriores.

Estas tres dimensiones son inherentes a las fuentes públicas disponibles. Su
efecto se mitiga parcialmente al trabajar con promedios institucionales sobre el
período completo, pero no se elimina. Por ello, los resultados del perfilamiento
deben entenderse como una caracterización del período 2019--2024 en su conjunto y
no como una serie temporal de perfiles anuales, con un margen de error en los
indicadores de complejidad mayor que el de las variables de volumen y estancia.

La tercera limitación es de alcance, dado que la base no contiene información sobre
dotación de personal, infraestructura física ni desenlaces clínicos post-alta.
En consecuencia, el perfilamiento describe lo que los hospitales registran durante
la hospitalización y no lo que producen en términos de salud de la población
atendida. Esta restricción proviene del diseño original de la base, que no fue concebida
para capturar resultados clínicos, y no puede resolverse mediante ningún proceso
de limpieza o transformación de los datos, dado que la información simplemente no
está registrada. Sin embargo, esta limitación es coherente con el objetivo del
estudio, que busca caracterizar el perfil funcional de los establecimientos a
partir de su actividad clínica registrada y no evaluar su desempeño en resultados.


\clearpage
\begin{table}[tp]
\centering
\caption{Principales variables del registro GRD y sus tipos de dato.}
\label{tab:variables-grd}
\begin{tabular}{p{4.2cm} p{7.0cm} p{1.8cm}}
\toprule
\textbf{Variable} & \textbf{Descripción} & \textbf{Tipo} \\
\midrule
\multicolumn{3}{l}{\textit{Identificación del establecimiento y paciente}} \\
\addlinespace
\texttt{COD\_HOSPITAL}        & Código del establecimiento tratante           & Código \\
\texttt{CIP\_ENCRIPTADO}      & Identificador encriptado del paciente (hasta 2023) & Código \\
\texttt{ID\_BENEFICIARIO}     & Identificador del beneficiario (desde 2024)   & Código \\
\addlinespace
\multicolumn{3}{l}{\textit{Datos del paciente}} \\
\addlinespace
\texttt{SEXO}                 & Sexo declarado del paciente                   & Texto \\
\texttt{FECHA\_NACIMIENTO}    & Fecha de nacimiento                           & Fecha \\
\texttt{PREVISION}            & Sistema previsional de salud                  & Texto \\
\addlinespace
\multicolumn{3}{l}{\textit{Información hospitalaria}} \\
\addlinespace
\texttt{SERVICIO\_SALUD}      & Servicio de Salud al que pertenece el hospital & Texto \\
\texttt{TIPO\_PROCEDENCIA}    & Origen de la derivación al establecimiento    & Texto \\
\texttt{TIPO\_INGRESO}        & Modalidad de ingreso (urgencia, programado, obstétrico, etc.) & Texto \\
\texttt{TIPO\_ACTIVIDAD}      & Tipo de episodio (hospitalización o CMA)      & Texto \\
\texttt{FECHA\_INGRESO}       & Fecha de ingreso al establecimiento           & Fecha \\
\texttt{FECHA\_ALTA}          & Fecha de egreso del establecimiento           & Fecha \\
\texttt{TIPOALTA}             & Condición de alta (domicilio, traslado, fallecido, etc.) & Texto \\
\texttt{USOSPABELLON}         & Número de usos de pabellón quirúrgico         & Numérico \\
\addlinespace
\multicolumn{3}{l}{\textit{Clasificación GRD}} \\
\addlinespace
\texttt{IR\_29301\_COD\_GRD}  & Código del Grupo Relacionado por Diagnóstico  & Código \\
\texttt{IR\_29301\_PESO}      & Peso relativo del GRD asignado                & Numérico \\
\texttt{IR\_29301\_SEVERIDAD} & Índice de severidad del episodio              & Entero \\
\texttt{IR\_29301\_MORTALIDAD}& Índice de mortalidad del episodio             & Entero \\
\addlinespace
\multicolumn{3}{l}{\textit{Codificación clínica}} \\
\addlinespace
\texttt{DIAGNOSTICO1}--\texttt{35}   & Diagnósticos CIE-10 (principal y comorbilidades) & Código \\
\texttt{PROCEDIMIENTO1}--\texttt{30} & Procedimientos CIE-9-MC                       & Código \\
\bottomrule
\end{tabular}
\end{table}

Los campos identificados como \emph{código} corresponden a etiquetas
categóricas, no a magnitudes: aunque en varios casos se expresan con dígitos,
no admiten operaciones aritméticas y el pipeline los lee como texto para evitar
la pérdida de ceros a la izquierda y de sufijos alfabéticos, presentes por
ejemplo en las secciones de la CIE-9-MC.

\section{Preparación de los datos}
\label{sec:etl}

La etapa de preparación transforma los seis archivos anuales heterogéneos en 
una única base depurada y estandarizada a nivel de registro. El proceso se articula en cinco pasos secuenciales: estandarización y limpieza,
clasificación de la actividad, aplicación de los criterios de elegibilidad,
construcción de variables e ingeniería de características, y
preprocesamiento y normalización.


\subsection{Estandarización y limpieza}
\label{sec:estandarizacion}

Cada archivo anual se lee con su codificación específica según lo indicado en la
Tabla~\ref{tab:heterogeneidad}, y se homologan los nombres de columnas, los
identificadores de establecimiento y los formatos de fecha. A partir de las
fechas de ingreso y alta se derivan dos variables fundamentales para el análisis:
la estancia hospitalaria, expresada en días, y la edad del paciente al momento
del ingreso, calculada desde la fecha de nacimiento. El identificador de paciente, en cambio, no forma parte de las columnas que el
pipeline solicita al leer los archivos: como se detalla en la
Sección~\ref{sec:comprension-negocio}, ninguna variable analítica se construye a
partir de él, de modo que el cambio de nombre de \texttt{CIP\_ENCRIPTADO} a
\texttt{ID\_BENEFICIARIO} en 2024 no exige homologación alguna ni una
equivalencia demostrada entre ambos campos.

En esta fase además se descartan los registros con \texttt{COD\_HOSPITAL} vacío y
aquellos cuya variable \texttt{TIPO\_ACTIVIDAD} no corresponde a ninguna de las
modalidades reconocidas por el estándar GRD. De los 5{,}808{,}536 registros
originales, 21 fueron eliminados por estas razones, todos ellos del año 2019 y
todos por presentar valores no estándar en \texttt{TIPO\_ACTIVIDAD}:
específicamente las categorías \texttt{DESCONOCIDO} (18 registros) y
\texttt{NO IDENTIFICADO} (3 registros), que no pueden asignarse a ninguna
modalidad de análisis. El criterio de \texttt{COD\_HOSPITAL} vacío se mantiene
como control de integridad, pero no descartó ningún registro: los seis archivos
identifican el establecimiento tratante en la totalidad de sus filas. Cabe
destacar que tampoco fue eliminado ningún registro por estancia negativa ni por
fechas inválidas, dado que ambas condiciones se trataron mediante truncamiento y
conversión a nulo respectivamente, conservando el registro en la base. El
resultado es una base depurada de 5{,}808{,}515 registros, lo que confirma que
los archivos GRD presentan alta consistencia interna. Los seis archivos anuales depurados se consolidan en
un único archivo en formato \texttt{parquet}, que permite su consulta eficiente
sin necesidad de cargar el conjunto completo en memoria.

\subsection{Clasificación de la actividad}
\label{sec:clasificacion-actividad}

El sistema GRD chileno registra episodios de hospitalización bajo cuatro tipos de
actividad. La hospitalización convencional corresponde a ingresos con estadía
nocturna en cama asignada, ya sea de origen programado o electivo, donde el
paciente permanece al menos una noche en el establecimiento. La hospitalización
en urgencia es un subconjunto que identifica específicamente los episodios
ingresados por la vía de urgencia que derivan en estadía nocturna. La
hospitalización diurna, en cambio, agrupa atenciones que requieren cama y
procedimientos durante el día pero que no implican estadía nocturna, como
quimioterapia, transfusiones o procedimientos de corta duración. La cirugía mayor
ambulatoria (CMA) agrupa procedimientos quirúrgicos o de alto costo que se
resuelven sin estadía nocturna, con alta el mismo día de la
intervención~\citep{TableauGRD}.

La distinción entre hospitalización convencional y CMA es relevante para el
perfilamiento porque ambas modalidades responden a lógicas funcionales distintas.
La CMA se asocia a procedimientos electivos de baja complejidad relativa y
estancia nula o mínima, mientras que la hospitalización convencional concentra los
episodios de mayor duración, severidad y consumo de recursos. Incorporar ambas en
los mismos indicadores de estancia y peso GRD distorsionaría los perfiles, por lo
que los indicadores se calculan por separado según modalidad. La distribución
de los registros por tipo de actividad se presenta en la
Tabla~\ref{tab:actividad}.


\begin{table}[htbp]
\centering
\caption{Distribución de registros por tipo de actividad tras la limpieza.}
\label{tab:actividad}
\begin{tabular}{lrr}
\toprule
\textbf{Tipo de actividad} & \textbf{Registros} & \textbf{\%} \\
\midrule
Hospitalización convencional     & 4{,}800{,}768 & 82{,}6 \\
Cirugía mayor ambulatoria (CMA)  &   889{,}868   & 15{,}3 \\
Hospitalización en urgencia      &    62{,}551   &  1{,}1 \\
Hospitalización diurna           &    55{,}328   &  1{,}0 \\
\midrule
\textbf{Total}                   & \textbf{5{,}808{,}515} & \textbf{100{,}0} \\
\bottomrule
\end{tabular}
\end{table}

\subsection{Criterios de elegibilidad}
\label{sec:elegibilidad}

La base depurada abarca 72 establecimientos con datos GRD en el período. Para
garantizar que los indicadores agregados sean comparables entre hospitales, se
aplica un criterio de inclusión: presencia en al menos tres de los seis años
del período. Este umbral responde a que los indicadores institucionales como
la estancia media, la proporción de CMA o la diversidad diagnóstica son
estimaciones estadísticas que requieren varios años de observación para no
depender de las particularidades de un único año. Un establecimiento con
datos en uno o dos años podría presentar valores atípicos debidos a factores
circunstanciales como cierres temporales, cambios de gestión o
incorporación reciente al sistema de información, que no reflejan su perfil
funcional habitual.

Una primera observación empírica reduce la relevancia práctica de elegir un
punto de corte exacto: la distribución real de años presentes entre los 72
establecimientos de la base depurada es marcadamente bimodal. Como muestra la
Tabla~\ref{tab:distribucion-anios}, 4 hospitales reportan datos en un único
año y 3 en dos años, todos ellos incorporados al sistema de información
recién en 2023 o 2024, mientras que los 65 hospitales restantes reportan
los seis años completos del período. Ningún establecimiento del
universo real tiene exactamente 3, 4 o 5 años de datos. En consecuencia,
cualquier umbral situado en ese rango (3, 4 o 5) produce exactamente el mismo
universo final de 65 hospitales elegibles: la elección del valor específico
``3'' no es, en los hechos, una decisión consecuente sobre qué establecimientos
se incluyen o excluyen, sino que opera sobre un vacío natural de los datos.

\begin{table}[htbp]
\centering
\caption{Distribución de años con datos GRD entre los 72 establecimientos de la base depurada.}
\label{tab:distribucion-anios}
\begin{tabular}{c c l}
\toprule
\textbf{Años presentes} & \textbf{N.º de hospitales} & \textbf{Interpretación} \\
\midrule
1 & 4  & Incorporación reciente (2023--2024) \\
2 & 3  & Incorporación reciente (2023--2024) \\
3--5 & 0  & (vacío en los datos reales) \\
6 & 65 & Cobertura completa del período \\
\bottomrule
\end{tabular}
\end{table}

Para cuantificar el segundo aspecto de la observación, se realizó
una simulación de error por submuestreo sobre los 65 hospitales con cobertura
completa. Tomando el valor de cada indicador calculado con los seis años como
referencia (``verdad''), se recalculó el mismo indicador usando todas las
combinaciones posibles de $k \in \{1, 2, 3, 4, 5\}$ años (o una muestra
aleatoria de 200 combinaciones cuando el total de combinaciones superaba ese
valor) y se midió el error relativo absoluto resultante, con un intervalo de
confianza del 95\,\% obtenido por bootstrap sobre la mediana. La
Tabla~\ref{tab:simulacion-anios} resume el error relativo mediano para las
seis variables tradicionales del pipeline.

\begin{table}[htbp]
\centering
\caption{Error relativo mediano (\%) al estimar cada indicador institucional con $k$ años en lugar de los 6 años completos, con intervalo de confianza 95\,\% (\textit{bootstrap}).}
\label{tab:simulacion-anios}
\begin{tabular}{c c c c c c c}
\toprule
\textbf{$k$ años} & \textbf{Egresos/año} & \textbf{Estancia} & \textbf{Peso GRD} & \textbf{Severidad} & \textbf{Mortalidad} & \textbf{Tasa CMA} \\
\midrule
1 & 9{,}3\,\% & 5{,}8\,\% & 4{,}2\,\% & 4{,}6\,\% & 3{,}8\,\% & 21{,}3\,\% \\
2 & 6{,}0\,\% & 4{,}5\,\% & 4{,}9\,\% & 3{,}8\,\% & 3{,}2\,\% & 10{,}9\,\% \\
\textbf{3} & \textbf{4{,}5\,\%} & \textbf{3{,}3\,\%} & \textbf{3{,}6\,\%} & \textbf{2{,}8\,\%} & \textbf{2{,}4\,\%} & \textbf{8{,}0\,\%} \\
4 & 3{,}0\,\% & 2{,}2\,\% & 2{,}4\,\% & 1{,}9\,\% & 1{,}6\,\% & 5{,}3\,\% \\
5 & 1{,}9\,\% & 1{,}1\,\% & 0{,}8\,\% & 0{,}9\,\% & 0{,}7\,\% & 4{,}2\,\% \\
\bottomrule
\end{tabular}
\end{table}

El error decrece de forma monótona con $k$ para las seis variables, pero el patrón de la reducción marginal es lo que sustenta el
umbral adoptado: el error promedio (sobre las seis variables) cae de
$8{,}2\,\%$ con un año a $4{,}1\,\%$ con tres años (una reducción de
$4{,}1$ puntos porcentuales) mientras que agregar un cuarto y un quinto año
reduce el error en $1{,}4$ y $1{,}2$ puntos porcentuales adicionales,
respectivamente. Es decir, el mayor beneficio marginal de sumar años se
obtiene entre $k=1$ y $k=3$. A partir de tres años, cada año adicional
reduce el error a una tasa aproximadamente constante y sustancialmente menor
que en el tramo inicial. Este patrón de rendimientos decrecientes, y no una
referencia externa, es lo que justifica fijar el umbral en tres años: es el
primer punto en el que la reducción del error empieza a estabilizarse. La
única excepción es la tasa de CMA, cuyo error relativo permanece
comparativamente alto incluso con tres años ($8{,}0\,\%$), lo esperable dado
que es una proporción calculada sobre un subconjunto minoritario de egresos
(los de modalidad CMA) y por tanto más sensible a la variación año a año que
los indicadores calculados sobre la totalidad de los egresos de
hospitalización.

Con esta evidencia, el umbral de tres años se sostiene como un punto de
compromiso cuantificado no arbitrario, aunque tampoco derivado de una regla
de decisión formal única entre la reducción de error que ofrece cada año
adicional de datos y el tamaño del universo de análisis disponible. El
detalle completo de la simulación (código, valores por hospital y por
combinación de años) se encuentra en
los artefactos de análisis reproducibles y la documentación de resultados del repositorio del
proyecto.

La aplicación de este criterio reduce el conjunto a 65 establecimientos, que
constituyen la población de análisis sobre la cual se construye la matriz
institucional descrita en la sección siguiente. El efecto de este criterio sobre
la partición final, y no solo sobre el error de los indicadores, se evalúa
mediante el análisis de sensibilidad definido en la
Sección~\ref{sec:sensibilidad-elegibilidad}.

La Tabla~\ref{tab:descartados} detalla los siete establecimientos excluidos, los
años en que reportaron datos y el total de egresos registrados en el período.
\begin{table}[htbp]
\centering
\caption{Establecimientos excluidos por no cumplir el criterio de presencia temporal.}
\label{tab:descartados}
\begin{tabular}{p{6.2cm} p{3.8cm} p{1.8cm} r}
\toprule
\textbf{Establecimiento} & \textbf{Servicio de Salud} & \textbf{Años con datos} & \textbf{Egresos} \\
\midrule
Instituto Traumatológico Dr. T. Gebauer & Metropolitano Suroriente & 2023, 2024 & 6{,}206 \\
Hospital de Constitución                & Del Maule               & 2024       & 2{,}796 \\
Hospital San Juan de Dios (Cauquenes)   & Del Maule               & 2024       & 4{,}461 \\
Hospital de Lota                        & Concepción              & 2023, 2024 & 7{,}706 \\
Hospital de Tomé                        & Talcahuano              & 2023, 2024 & 6{,}365 \\
Hospital Penco-Lirquén                  & Concepción              & 2024       & 2{,}751 \\
Complejo Asistencial Padre Las Casas    & Araucanía Sur           & 2024       & 6{,}389 \\
\bottomrule
\end{tabular}
\end{table}

\subsection{Construcción de variables e ingeniería de características}
\label{sec:variables}

A partir de los registros depurados se construye un conjunto de aproximadamente
35 variables agregadas a nivel institucional, organizadas en cuatro grupos
conceptuales. Una decisión metodológica central es la exclusión de recursos de infraestructura y equipamiento tales como la dotación de camas o ambulancias, dado que el objetivo es un perfilamiento funcional y no
estructural, todas las variables se derivan exclusivamente de la actividad clínica
registrada, de modo que los grupos resultantes reflejen lo que los hospitales
hacen y no la capacidad física que poseen.

Una segunda decisión de diseño es utilizar la totalidad de los campos de
codificación clínica disponibles. Aunque el registro GRD admite hasta 35
diagnósticos CIE-10 y hasta 30 procedimientos CIE-9-MC por egreso
(Tabla~\ref{tab:variables-grd}), el pipeline detecta dinámicamente y procesa
\texttt{DIAGNOSTICO1}--\texttt{DIAGNOSTICO35} y
\texttt{PROCEDIMIENTO1}--\texttt{PROCEDIMIENTO30}. Para controlar el costo de
memoria de procesar 65 campos clínicos sobre 5{,}8 millones de registros, los
conteos se acumulan columna a columna en vez de materializar una tabla extensa
de todos los códigos. Así, el promedio de comorbilidades corresponde al número
medio de diagnósticos secundarios registrados entre
\texttt{DIAGNOSTICO2} y \texttt{DIAGNOSTICO35}, con rango teórico $[0,34]$,
y los vectores ponderados de capítulos CIE-10 y de secciones CIE-9-MC reflejan
la codificación completa del episodio.

Esta cobertura evita subrepresentar los egresos con mayor profundidad de
registro, pero no elimina la heterogeneidad de codificación entre
establecimientos. Por ello, las diferencias en comorbilidades deben interpretarse
como diferencias en carga clínica \emph{registrada}, potencialmente afectadas por
la exhaustividad de codificación de cada centro, y no como una medida clínica
pura. El análisis conserva los indicadores porque incorporan información
relevante para el perfilamiento relativo, y declara esta limitación al
interpretar los resultados.

\clearpage
La Tabla~\ref{tab:variables} sintetiza las variables construidas por grupo.

{\small
\setstretch{1.0}
\setlength{\LTleft}{0pt}
\setlength{\LTright}{\fill}
\begin{longtable}{>{\raggedright\arraybackslash}p{3.2cm} >{\raggedright\arraybackslash}p{8.5cm}}
\caption{Variables construidas a nivel institucional, agrupadas por dimensión.}
\label{tab:variables} \\
\toprule
\textbf{Grupo} & \textbf{Variables} \\
\midrule
\endfirsthead
\multicolumn{2}{l}{\small\textit{Tabla~\thetable\ (continuación)}} \\
\toprule
\textbf{Grupo} & \textbf{Variables} \\
\midrule
\endhead
\bottomrule
\endfoot
\raggedright Tradicionales de gestión &
Egresos por año, estancia media, estancia mediana, peso medio GRD, índice de
severidad medio, índice de mortalidad medio. \\
\addlinespace
\raggedright Diversidad diagnóstica &
Entropía de Shannon normalizada sobre la distribución de GRD, comorbilidades
promedio por egreso (diagnósticos secundarios 2 a 35, rango teórico $[0,34]$). \\
\addlinespace
\raggedright Cirugía mayor ambulatoria &
Tasa de CMA, peso medio de los GRD de CMA. \\
\addlinespace
\raggedright Perfil funcional agregado &
Distribución demográfica (\% pediátrico, \% geriátrico, \% femenino fértil), mezcla
por tipo de ingreso (\% urgencia, \% programado, \% referencia), tipo de alta (\%
domicilio, \% traslado, \% fallecido), uso de pabellón, actividad obstétrica (tasa
de partos, \% obstétrica, tasa de prematurez), dispersión de la estancia
(coeficiente de variación). \\
\bottomrule
\end{longtable}
}

\textbf{Variables tradicionales de gestión:} Este grupo captura los indicadores de producción y complejidad habituales en la
gestión hospitalaria. El peso medio GRD es especialmente relevante porque resume
en un único valor la intensidad de recursos esperada para la casuística promedio
del establecimiento, mientras que los índices de severidad y mortalidad reflejan
la gravedad clínica de la población atendida.

\textbf{Variables de diversidad diagnóstica:} Para cuantificar la amplitud del espectro clínico de cada hospital se emplea la
entropía de Shannon normalizada sobre la distribución de GRD. Dada la
distribución de frecuencias relativas $p_i$ de los $G$ grupos GRD distintos
presentes en un hospital, el indicador se define como:

\begin{equation}
H_{\text{norm}} = \frac{-\sum_{i=1}^{G} p_i \, \log_2(p_i)}{\log_2 G}
\label{eq:shannon}
\end{equation}

La división por $\log_2 G$, que es el valor máximo que puede alcanzar la
entropía cuando la distribución es uniforme, acota el indicador al intervalo
$[0, 1]$ y lo hace comparable entre hospitales con distinto número de GRD
distintos; por convención se asigna el valor 0 cuando $G \leq 1$ y se excluyen
los términos con $p_i = 0$. Un hospital con entropía normalizada alta atiende una
gran variedad de condiciones con frecuencias relativamente uniformes (perfil
generalista), mientras que uno con entropía baja concentra su actividad en pocos
grupos (perfil especializado). Esta métrica complementa los indicadores de
complejidad al capturar la diversidad de la producción.

\textbf{Variables de cirugía mayor ambulatoria:} La tasa de CMA y el peso medio de sus GRD caracterizan la orientación ambulatoria del establecimiento, una combinación que resulta especialmente informativa, ya que una tasa alta con peso bajo señala un perfil orientado a procedimientos electivos simples, mientras que una tasa alta con peso elevado indica un programa de CMA de mayor complejidad técnica en el establecimiento. Cabe precisar que, de estas dos variables, solo la tasa de CMA integra la matriz sobre la que opera el agrupamiento: el peso medio de CMA no está definido para los establecimientos sin actividad de esta modalidad, por lo que se conserva con fines exclusivamente descriptivos y se excluye del modelado, como se detalla en la Sección~\ref{sec:matriz-final}.

\textbf{Variables de perfil funcional agregado:}
Este grupo, el más numeroso, describe la composición de la actividad clínica
mediante proporciones calculadas exclusivamente sobre los egresos de
hospitalización (excluyendo CMA). La separación no responde a ausencia de información: los campos demográficos,
de ingreso, alta, procedencia y pabellón también están disponibles en los
episodios CMA. Responde a que la CMA es una modalidad ambulatoria: en los datos, el $93{,}3\,\%$ de sus episodios registra estancia igual a cero. Mezclarla con hospitalización alteraría el significado y los denominadores de indicadores definidos para hospitalización. Por ejemplo, un
hospital con mayor proporción de CMA reduciría mecánicamente sus tasas de
estancia prolongada o de ingreso obstétrico, aun sin modificar su perfil de
hospitalización. La orientación ambulatoria no se omite:
se incorpora separadamente mediante la tasa de CMA, mientras que el peso medio
de los GRD de CMA se conserva como variable descriptiva. Cada variable de este
bloque se define operacionalmente como la fracción de egresos hospitalarios que
cumple una condición, lo que las hace comparables entre hospitales de distinto
volumen. La Tabla~\ref{tab:definiciones} detalla sus definiciones operacionales.

\begin{table}[htbp]
\centering
\caption{Definiciones operacionales de las variables de perfil funcional
agregado.}
\label{tab:definiciones}
\begin{tabular}{p{5.2cm} p{7.0cm}}
\toprule
\textbf{Variable} & \textbf{Definición operacional} \\
\midrule
\% pediátrico & Proporción de egresos con edad al ingreso menor a 15 años. \\
\% geriátrico & Proporción de egresos con edad al ingreso mayor o igual a 65
años. \\
\% femenino fértil & Proporción de egresos de mujeres entre 15 y 49 años. \\
Edad mediana & Mediana de la edad al ingreso de los egresos. \\
\% urgencia / programada / obstétrica & Proporción de egresos según tipo de
ingreso. \\
\% alta domicilio / fallecido / traslado & Proporción de egresos según tipo de
alta. \\
\% origen emergencia / referencia & Proporción según procedencia del paciente
(servicio de urgencia propio vs.\ derivación de otro establecimiento). \\
\% uso de pabellón & Proporción de egresos con al menos un uso de pabellón
quirúrgico. \\
Tasa de partos & Proporción de egresos con condición de alta neonatal registrada.
\\
Tasa de prematurez & Proporción de partos con peso del recién nacido inferior a
2.500 g. \\
\% estancia larga & Proporción de egresos con estancia superior a 30 días. \\
CV de estancia & Coeficiente de variación de la duración de estancia. \\
\bottomrule
\end{tabular}
\end{table}

Estas variables permiten distinguir perfiles funcionales específicos. Una alta
tasa de partos y proporción obstétrica caracteriza a establecimientos con
maternidad activa, mientras que una alta proporción pediátrica identifica
hospitales de niños. De forma similar, una elevada proporción de origen por
referencia señala centros que operan como referentes de su red territorial. A
estas dimensiones se suman el coeficiente de variación de la estancia y la
proporción de estancias largas, que en conjunto capturan la heterogeneidad
interna del proceso de hospitalización y complementan los perfiles anteriores
con información sobre la variabilidad de la gestión clínica dentro de cada
establecimiento.


\subsubsection{Reducción dimensional de la casuística clínica}
\label{sec:pca}

La casuística clínica de cada hospital se resume mediante tres vectores de
proporciones construidos sobre sus egresos. El primero corresponde a la
distribución por los 22 capítulos de la CIE-10. El segundo recoge la
distribución por las 18 secciones de procedimientos definidas en la tabla
maestra CIE-9-MC empleada, es decir, las 17 secciones numeradas de 00 a 16 y la
sección 03A. El tercero describe la concentración del establecimiento en los
20 GRD más frecuentes a nivel nacional. Cada vector incluye una categoría
residual. En los dos primeros casos reúne los códigos que no admiten mapeo a un
capítulo o una sección de la tabla maestra; en el tercero agrupa los GRD que
quedan fuera del Top-20 nacional. Estas categorías permiten que cada vector
forme una distribución completa cuya suma es igual a 1.

Al concatenar los tres vectores se obtiene una matriz de $65 \times 63$
columnas: 23 de capítulos CIE-10, 19 de secciones CIE-9-MC y 21 de GRD, con la
categoría residual incluida en cada caso. Usar directamente esta matriz para
solo 65 hospitales introduce un problema de alta dimensionalidad relativa. En
ese escenario, las distancias entre observaciones pierden capacidad
discriminatoria, el agrupamiento se vuelve inestable y la interpretación se
dificulta, como se argumentó en la Sección~\ref{sec:reduccion}. Por ello, antes
de integrar la casuística con el resto de las variables se aplica PCA para
reducir las 63 columnas a un conjunto compacto de componentes principales.

Estos tres vectores son, en sentido estricto, datos composicionales, pues cada
fila es una distribución de proporciones que suma 1. Aplicar PCA y, más
adelante, distancia euclidiana directamente sobre estas proporciones puede
introducir el problema de la \emph{clausura composicional} (\textit{closure
problem}). Como las partes deben sumar una constante, el incremento de una
proporción obliga necesariamente a la disminución de otras. Esto puede generar
correlaciones y distancias artificiales que no reflejan diferencias reales en la
actividad clínica. La Sección~\ref{sec:sensibilidad-composicional} define el
análisis que evalúa empíricamente la magnitud de este efecto sobre los
resultados del trabajo.

\subsubsection{Análisis de componentes principales}
\label{sec:pca-detalle}

Sobre la matriz concatenada de casuística, previa estandarización de las
proporciones, se aplica PCA. Dado que el vector de capítulos CIE-10 admite dos
construcciones, se evaluaron ambas variantes de forma independiente. La variante
\emph{principal} considera únicamente el diagnóstico principal
(\texttt{DIAGNOSTICO1}) de cada egreso, de modo que cada egreso aporta una
unidad al capítulo de su causa de hospitalización. La variante \emph{ponderada}
incorpora además los diagnósticos secundarios registrados, asignando peso
$1{,}0$ al diagnóstico principal y $0{,}5$ a cada secundario, con el propósito
de reflejar las comorbilidades del episodio sin que estas lleguen a dominar
sobre el motivo principal de ingreso. La ponderación es, por tanto, por posición
del diagnóstico dentro del episodio y no por el peso relativo del GRD asignado.
La comparación entre ambas variantes se resolvió mediante un criterio interno al
propio PCA:
para cada variante se ejecutó el PCA con el número de componentes fijado en
$K=8$ (justificado más adelante en esta sección) y se comparó la varianza
acumulada explicada, sin recurrir a ninguna métrica de clustering posterior.
La variante ponderada explica el 72{,}3\,\% de la varianza acumulada frente al
71{,}8\,\% de la variante basada solo en el diagnóstico principal, por lo que fue seleccionada
para el pipeline definitivo. Este criterio de selección --varianza explicada
por el propio PCA-- es deliberadamente independiente del clustering que se
describe en las secciones siguientes, para evitar que la elección de una
decisión de preprocesamiento quede determinada por la misma métrica con la
que luego se evalúa la solución final. Esta misma comparación se repite bajo una
representación composicional alternativa (CLR) como parte del análisis de
sensibilidad definido en la Sección~\ref{sec:sensibilidad-composicional}.

El número de componentes a retener se fija en $K=8$, decisión que se
verificó empíricamente mediante un barrido de $K \in \{8, 9, 10, 11\}$
evaluando el efecto de cada valor sobre las métricas de validación interna
del clustering final (Tabla~\ref{tab:barrido-k-pca}). Con la variante
ponderada, las 8 componentes explican el 72{,}3\,\% de la varianza acumulada
de las 63 columnas de casuística. Incorporar componentes adicionales
incrementa la varianza retenida, pero no mejora la calidad de la partición
resultante. Por el contrario, al aumentar $K$ de 8 a 11 --el número mínimo
que acumularía el 80\,\% de la varianza-- el coeficiente de Silhouette
desciende de $0{,}401$ a $0{,}380$, el índice de Calinski-Harabasz cae de
$22{,}4$ a $19{,}7$ y el índice de Davies-Bouldin empeora de $0{,}990$ a
$1{,}076$. El valor $K=9$ obtiene el mejor Silhouette del barrido
($0{,}422$), pero fragmenta los grupos especializados de forma menos
interpretable que $K=8$: reintroduce un grupo de apenas dos establecimientos
(tamaños $3$, $54$, $2$, $6$ frente a $3$, $54$, $5$, $3$), precisamente la
situación en que las pruebas de significancia basadas en aproximaciones
asintóticas pierden validez (Sección~\ref{sec:nivel2-metodo}). Se retienen,
por tanto, las 8 componentes por ser el valor que ofrece el mejor balance
entre las tres métricas de validación interna y la interpretabilidad de los
grupos resultantes, y no un porcentaje de varianza explicada fijado de
antemano: agregar componentes de casuística más allá de este punto introduce
ruido de alta dimensionalidad relativa (Sección~\ref{sec:pca}) que degrada la
cohesión y separación del clustering en mayor medida que lo que aporta en
varianza retenida.

\begin{table}[htbp]
\centering
\caption{Efecto del número de componentes PCA de casuística ($K$) sobre las métricas de validación interna del clustering Ward $K=4$.}
\label{tab:barrido-k-pca}
\begin{tabular}{c c c c c l}
\toprule
\textbf{$K$ (componentes)} & \textbf{Varianza acum.} & \textbf{Silhouette} & \textbf{CH} & \textbf{DB} & \textbf{Tamaños} \\
\midrule
\bfseries 8 & \bfseries 72{,}3\,\% & 0{,}401 & \bfseries 22{,}4 & \bfseries 0{,}990 & 3, 54, 5, 3 \\
9 & 75{,}2\,\% & \bfseries 0{,}422 & 21{,}9 & 0{,}991 & 3, 54, 2, 6 \\
10 & 77{,}7\,\% & 0{,}384 & 20{,}4 & 1{,}045 & 3, 54, 5, 3 \\
11 & 80{,}0\,\% & 0{,}380 & 19{,}7 & 1{,}076 & 3, 54, 5, 3 \\
\bottomrule
\end{tabular}
\par\smallskip
\begin{minipage}{\textwidth}
\footnotesize\raggedright
\textit{Nota.} CH corresponde al índice de Calinski--Harabasz (un valor mayor indica mejor separación entre grupos); DB corresponde al índice de Davies--Bouldin (un valor menor indica grupos más compactos y separados).
\end{minipage}
\end{table}

Cada componente admite una lectura interpretable a partir de sus cargas. La
Tabla~\ref{tab:cargas-pca} reporta, para cada una de las 8 componentes, las
tres variables originales con mayor magnitud de carga y su signo, lo que
permite sostener empíricamente su interpretación clínica en lugar de
describirla solo en términos generales.

{\footnotesize
\setstretch{1.0}
\setlength{\LTleft}{0pt}
\setlength{\LTright}{\fill}
\begin{longtable}{@{}>{\raggedright\arraybackslash}p{1.4cm}@{\hspace{0.25cm}}>{\raggedright\arraybackslash}p{3.45cm}@{\hspace{0.25cm}}>{\raggedright\arraybackslash}p{3.45cm}@{\hspace{0.25cm}}>{\raggedright\arraybackslash}p{3.45cm}@{}}
\caption{Tres cargas de mayor magnitud por componente PCA de casuística. El signo indica la dirección de la carga.}
\label{tab:cargas-pca} \\
\toprule
\textbf{Componente} & \textbf{Carga 1} & \textbf{Carga 2} & \textbf{Carga 3} \\
\midrule
\endfirsthead
\multicolumn{4}{l}{\footnotesize\textit{Tabla~\thetable\ (continuación)}} \\
\toprule
\textbf{Componente} & \textbf{Carga 1} & \textbf{Carga 2} & \textbf{Carga 3} \\
\midrule
\endhead
\bottomrule
\endfoot
\texttt{dim\_01} & $-$ Embarazo, parto y puerperio ($-0{,}24$) & $+$ Otros GRD fuera del Top-20 ($0{,}23$) & $-$ Procedimientos obstétricos ($-0{,}22$) \\
\addlinespace[3pt]
\texttt{dim\_02} & $-$ Enfermedades de la piel y tejido subcutáneo ($-0{,}27$) & $+$ Operaciones sobre el sistema nervioso ($0{,}24$) & $+$ Operaciones sobre el sistema endocrino ($0{,}23$) \\
\addlinespace[3pt]
\texttt{dim\_03} & $-$ Enfermedades del aparato circulatorio ($-0{,}32$) & $+$ Operaciones sobre el oído ($0{,}27$) & $+$ Malformaciones congénitas y anomalías cromosómicas ($0{,}25$) \\
\addlinespace[3pt]
\texttt{dim\_04} & $-$ PH procedimientos sobre extremidad superior ($-0{,}34$) & $-$ PH procedimientos sobre rodilla y extremidad inferior ($-0{,}31$) & $-$ Operaciones sobre el aparato musculoesquelético ($-0{,}31$) \\
\addlinespace[3pt]
\texttt{dim\_05} & $+$ Operaciones sobre el aparato respiratorio ($0{,}30$) & $-$ Enfermedades del ojo y sus anexos ($-0{,}27$) & $-$ Trastornos mentales y de comportamiento ($-0{,}26$) \\
\addlinespace[3pt]
\texttt{dim\_06} & $+$ Operaciones sobre el aparato digestivo ($0{,}37$) & $+$ Otros procedimientos diagnósticos y terapéuticos ($0{,}33$) & $+$ Operaciones sobre el sistema hemático y linfático ($0{,}32$) \\
\addlinespace[3pt]
\texttt{dim\_07} & $+$ PH colecistectomía laparoscópica ($0{,}36$) & $-$ MH trastornos del anteparto con complicaciones ($-0{,}33$) & $+$ PH cesárea ($0{,}26$) \\
\addlinespace[3pt]
\texttt{dim\_08} & $-$ Operaciones sobre el sistema endocrino ($-0{,}28$) & $+$ MH trastornos del anteparto con complicaciones ($0{,}26$) & $-$ Otros procedimientos diagnósticos y terapéuticos ($-0{,}22$) \\
\bottomrule
\end{longtable}
}

Estas componentes, denotadas \texttt{dim\_01} a \texttt{dim\_08}, reemplazan
las 63 columnas originales de casuística, condensando el perfil
clínico-quirúrgico de cada hospital en un espacio compacto interpretable. Las
cargas completas de las 8 componentes sobre las 63 variables originales se
reportan en el archivo \texttt{reports/tables/reduccion\_dimensional.csv} del
repositorio del proyecto.

\subsubsection{Matriz final de variables}
\label{sec:matriz-final}
La integración de los cuatro grupos de
variables institucionales directas junto con las 8 componentes principales de
casuística conforma la matriz de trabajo definitiva, denominada
\textit{Hospital Matrix Integrada}:

\begin{equation}
  \underbrace{27}_{\text{variables directas}} + \underbrace{8}_{\text{componentes PCA}} = \mathbf{35 \text{ columnas}}
\end{equation}

Las 27 variables directas se distribuyen en 6 tradicionales de gestión, 2 de
diversidad diagnóstica, 1 de cirugía mayor ambulatoria (la tasa de CMA) y 18 de
perfil funcional agregado. La segunda variable de CMA construida en la
Sección~\ref{sec:variables}, el peso medio de los GRD de esta modalidad, queda
fuera de la matriz de modelado porque no está definida para los establecimientos
que no registran actividad ambulatoria en el período: mantenerla habría exigido
imputar un valor arbitrario en esos casos o eliminar hospitales del universo de
análisis. Se conserva, por tanto, únicamente como variable descriptiva en la
caracterización de los perfiles.

Esta matriz tiene dimensiones de $65 \times 35$, donde cada fila representa un
hospital y cada columna una característica derivada exclusivamente de su actividad
clínica. Las 63 columnas originales de casuística no forman parte de la matriz
final, ya que fueron sustituidas íntegramente por las 8 componentes PCA. Sobre esta
matriz opera todo el pipeline de agrupamiento, validación e interpretación
descrito en las secciones siguientes.

\subsection{Preprocesamiento y normalización}
\label{sec:preprocesamiento}

Antes del agrupamiento, las variables se someten a transformaciones
secuenciales orientadas a reducir la influencia de distribuciones con cola
pronunciada y a hacer comparables las escalas de los indicadores en el cálculo
de distancias euclidianas.

\textbf{Transformación logarítmica:} Se aplicó la transformación $\log(1+x)$
a las variables \textit{egresos por año} y \textit{pabellones promedio}, ambas
con asimetría positiva. El volumen anual de egresos varía entre 2.663 y 39.066,
con una asimetría de $0{,}94$, lo que refleja la presencia de algunos
establecimientos de gran tamaño relativo. El promedio de usos de pabellón
presenta la mayor asimetría positiva del conjunto ($2{,}01$), concentrando a la
mayoría de los hospitales en valores bajos y a unos pocos en valores elevados.
Para ambas variables se utilizó:

\begin{equation}
x' = \log(1 + x)
\label{eq:log1p}
\end{equation}

La transformación reduce la influencia de los valores altos sobre las distancias
entre hospitales y evita la indeterminación de $\log(x)$ cuando $x = 0$. La edad
mediana no se transformó logarítmicamente: presenta asimetría negativa
($-1{,}38$), y $\log(1+x)$ está orientada a comprimir colas derechas, por lo que
su aplicación habría acentuado ese patrón. Las variables restantes no se
transformaron logarítmicamente, pues no presentaban colas derechas de magnitud
comparable; sus diferencias de escala se abordan mediante el escalado robusto.

\textbf{Escalado robusto:} Posteriormente, todas las variables de la matriz,
incluidas las transformadas, se estandarizaron mediante un escalador robusto que
las centra en la mediana y las escala por el rango intercuartílico (IQR):

\begin{equation}
z = \frac{x' - \operatorname{mediana}(x')}{\operatorname{IQR}(x')}
\label{eq:robust}
\end{equation}

En las variables que no requirieron transformación logarítmica, $x'$ corresponde
a su valor original. La elección de estos estadísticos frente a los clásicos
basados en media y desviación estándar responde a la presencia de valores
extremos legítimos en la red. El Hospital Sótero del Río, por ejemplo, registra
el mayor volumen de egresos del sistema con 39.066 egresos anuales, cifra que no
constituye un error de registro sino una realidad operativa. Un escalador clásico
sería arrastrado por este extremo, comprimiendo artificialmente la escala del
resto de los hospitales. El escalador robusto atenúa este efecto al basar la
estandarización en la mediana y el IQR, que son resistentes a valores extremos,
y permite conservar estos establecimientos en el análisis.

Como verificación de sensibilidad, se comparó la partición Ward con $K=4$
obtenida con y sin transformación logarítmica de la edad mediana. Ambas
configuraciones produjeron la misma partición de los 65 hospitales
(ARI $= 1{,}00$), por lo que la corrección mejora la coherencia del
preprocesamiento sin alterar los perfiles funcionales reportados.

\subsubsection{Diagnóstico y sensibilidad de bloques de variables}
\label{sec:bloques-variables}

Dado que la distancia euclidiana acumula la contribución de cada variable, se
cuantificó la fracción de la suma de distancias cuadradas entre todos los pares
de hospitales atribuible a cuatro macro-bloques de la matriz ya preprocesada:
gestión, diversidad y CMA (9 variables); perfil operativo no demográfico (11);
demanda demográfica-obstétrica (7); y componentes PCA de casuística (8). El
bloque demográfico-obstétrico comprende porcentaje pediátrico, geriátrico y
femenino fértil, edad mediana, porcentaje de ingreso obstétrico, proporción de
egresos neonatales y tasa de prematurez. Esta definición distingue la
composición de demanda de las variables operativas y no atribuye a priori las
componentes PCA a un bloque demográfico, pues cada componente resume una mezcla
de variables de casuística.

Como sensibilidad de ponderación, después del escalado robusto cada variable de
un macro-bloque se multiplica por $1/\sqrt{p_b}$, donde $p_b$ es el número de
variables del bloque. Esta transformación evita que un bloque tenga mayor peso
solo por contener más columnas. Además, se re-ejecuta Ward $K=4$ excluyendo las
siete variables demográfico-obstétricas. En ambos escenarios se comparan las
particiones con la solución principal mediante ARI y NMI, y se reportan
Silhouette, Calinski--Harabasz, Davies--Bouldin y tamaños de grupo. Los
resultados se presentan en la Sección~\ref{sec:res-bloques}.

% ============================================================================
% ============================================================================
\section{Modelado}
\label{sec:metodo-clustering}
% ============================================================================

Esta sección corresponde a la etapa de modelado de CRISP-DM y describe el
procedimiento de agrupamiento aplicado sobre la matriz institucional construida
en la Sección~\ref{sec:etl}. El diseño del procedimiento responde a dos
decisiones interdependientes: la elección del algoritmo y la organización en dos
niveles de análisis.


\subsection{Elección del algoritmo}
\label{sec:algoritmo}


Se aplica clustering aglomerativo jerárquico con el método de Ward sobre
la matriz preprocesada, utilizando la distancia euclidiana como métrica
de disimilitud entre observaciones.
La elección sobre alternativas particionales como
$k$-means responde a una propiedad del problema: con $n = 65$ establecimientos y
sin consenso institucional sobre cuántos perfiles funcionales deberían existir
en el sistema público chileno, fijar $k$ de antemano introduce un supuesto
arbitrario. El dendrograma del método jerárquico permite, en cambio, explorar
la estructura a distintos niveles de granularidad antes de comprometerse con una
partición particular, haciendo visible cómo se fusionan los hospitales y a qué
costo de cohesión.

El criterio de enlace de Ward minimiza, en cada fusión, el incremento de la suma
de cuadrados intra-grupo. Para dos clusters $A$ y $B$, el criterio es:

\begin{equation}
  \Delta(A, B) = \frac{|A|\,|B|}{|A| + |B|}
                 \,\lVert \mathbf{c}_A - \mathbf{c}_B \rVert^2
  \label{eq:ward}
\end{equation}

donde $|A|$ y $|B|$ son los tamaños de los clusters y $\mathbf{c}_A$,
$\mathbf{c}_B$ sus centroides. Este criterio produce grupos compactos y
relativamente equilibrados, propiedad deseable cuando el objetivo es maximizar
la homogeneidad funcional dentro de cada perfil.


\subsection{Motivación del enfoque en dos niveles}
\label{sec:motivacion-jerarquico}

La organización del procedimiento en dos niveles surge de la evidencia acumulada durante la exploración
algorítmica. Para establecer cuántos grupos existen en la red, se ejecutaron
ocho técnicas que determinan $k$ automáticamente sin que el analista lo fije:
HDBSCAN, Mean Shift, Affinity Propagation, clustering aglomerativo con criterio
\textit{largest gap}, estadístico Gap, MST-kNN~\citep{Villalobos2016},
KNN-graph con Leiden y GMM con criterio BIC. Las ocho convergen en la misma
lectura: la red presenta un núcleo numeroso de hospitales con perfiles similares
entre sí más un conjunto reducido de establecimientos que se distancian
marcadamente del resto.

Ante esta estructura, un agrupamiento global único resulta insatisfactorio en
ambos extremos del rango. Con $k$ bajo, el corte más fuerte se produce en
$k = 2$ (Silhouette $= 0{,}605$), que aísla un grupo pequeño del resto pero
deja al 95\,\% de la red como un bloque indiferenciado, sin valor para el
perfilamiento. Con $k$ alto, al explorar $k \in \{5, \ldots, 20\}$ con
K-means++, clustering aglomerativo Ward y GMM, el coeficiente de Silhouette
decrece de forma consistente en las tres técnicas a medida que $k$ crece, lo
que indica que el sistema no soporta una segmentación más granular sobre los
65 hospitales en conjunto.

Para resolver esta dualidad, el procedimiento se organiza en dos niveles
independientes:

\begin{itemize}
  \item \textbf{Nivel 1 — partición global}: opera sobre los 65 hospitales y
    separa los grupos funcionales diferenciados del núcleo general.
  \item \textbf{Nivel 2 — subdivisión del núcleo generalista}: opera exclusivamente
    sobre los hospitales asignados al grupo principal del Nivel~1 y explora
    la estructura interna de ese conjunto, que el Nivel~1 trata como un bloque
    homogéneo.
\end{itemize}

Los grupos concretos obtenidos, su tamaño y sus perfiles característicos se
presentan en el Capítulo~\ref{cap:resultados}.

\subsection{Selección del número de grupos}
\label{sec:seleccion-k}

El número de grupos se examina reportando conjuntamente tres métricas internas
para $k \in \{2, 3, \ldots, 8\}$: el coeficiente de Silhouette, el índice de
Calinski-Harabasz y el índice de Davies-Bouldin, definidos y justificados en la
Sección~\ref{sec:validacion}. Estas métricas no se emplean como una regla de
decisión que determine automáticamente un valor de $k$, porque los valores más
bajos de $k$ tienden a optimizar simultáneamente las tres métricas al concentrar
la mayor parte de la red en un único bloque y aislar solo a los establecimientos
más extremos, produciendo particiones con escasa utilidad para el perfilamiento
buscado. En sentido contrario, valores más altos de $k$ subdividen mecánicamente
grupos que en realidad son homogéneos, generando perfiles que se distinguen por
diferencias marginales en lugar de reflejar patrones funcionales genuinamente
distintos. Por este motivo, el valor de $k$ finalmente adoptado no corresponde
al óptimo matemático de ninguna métrica individual ni a un punto de convergencia
de las tres, sino a un compromiso interpretativo: se selecciona la partición de
menor $k$ que separa perfiles funcionalmente distinguibles según el conocimiento
de dominio, aun cuando particiones con $k$ menor obtengan mejores valores en las
tres métricas internas. Esta decisión se declara explícitamente como tal en el
Capítulo~\ref{cap:resultados}, evitando calificar al $k$ elegido como
``número óptimo de grupos''.

Para interpretar el coeficiente de Silhouette se toman como referencia los
umbrales propuestos por \citet{Rousseeuw1987}: valores superiores a
0{,}70 indican una estructura ``fuerte'', entre 0{,}50 y 0{,}70 una estructura
``razonable'', y por debajo de 0{,}50 una estructura ``débil''. Sin embargo,
estos umbrales fueron concebidos para datos de baja dimensionalidad con separación
nítida entre grupos. Con matrices de características reales como la empleada en
este trabajo, las distancias se distribuyen de forma más uniforme y el coeficiente
raramente alcanza los rangos considerados ``fuertes''. Para situar esta expectativa
en el contexto del problema, \citet{Giglio2021} evaluaron la
clasificación administrativa del MINSAL con el mismo coeficiente y obtuvieron
valores inferiores a 0{,}2 en todos los años analizados, con valores negativos en
algunos casos. Por este motivo, los umbrales se emplean como referencia
orientativa para comparar soluciones entre sí, no como criterio de aceptación o
rechazo. Las métricas para cada valor de $k$ y la justificación de la partición
adoptada se presentan en el Capítulo~\ref{cap:resultados}.

El mismo criterio multimétrico se aplica de forma independiente en el Nivel~2,
sobre el subconjunto de hospitales del grupo mayoritario identificado en el
Nivel~1.

\subsection{Caracterización post-hoc de los grupos}
\label{sec:distribucion}

Una vez obtenida la partición, cada grupo se caracteriza mediante un
procedimiento \textit{post-hoc} calculando las medianas de las variables
originales (en escala interpretable) por cluster y se identifican los atributos
en los que cada grupo se distingue del resto. A partir de este perfil de
medianas se asigna a cada grupo una etiqueta funcional descriptiva.

El aspecto metodológicamente relevante es que esta etiquetación se realiza
después del clustering y no se entrega ninguna etiqueta al
algoritmo. El agrupamiento opera exclusivamente sobre las variables de actividad
clínica. Este
orden garantiza que los perfiles describan patrones emergentes de los datos y no
categorías impuestas antes. Los grupos concretos obtenidos, su tamaño
y sus perfiles característicos se presentan y discuten en el
Capítulo~\ref{cap:resultados}.

\subsection{Subdivisión del núcleo generalista (Nivel 2)}
\label{sec:nivel2-metodo}

El grupo mayoritario identificado en el Nivel 1 concentra la mayor parte de los
65 hospitales del sistema y presenta heterogeneidad interna que el Nivel 1 no
discrimina al tratarlo como un bloque homogéneo. Para explorar esta estructura,
se aplica nuevamente clustering aglomerativo Ward sobre el subconjunto de
hospitales del grupo principal. El número de subgrupos se selecciona con el
mismo criterio multimétrico de la Sección~\ref{sec:seleccion-k}.

Dado que el análisis opera sobre un subconjunto de $n = 54$ hospitales, el
tamaño muestral efectivo del Nivel 2 es menor que el del Nivel 1, y los
resultados deben interpretarse como exploratorios: sirven para generar
hipótesis sobre matices funcionales dentro del grupo principal, como
gradientes de complejidad o especialización parcial, y no para establecer
fronteras funcionales con el mismo nivel de certeza que el Nivel 1. Esta
cautela es particularmente necesaria cuando alguno de los subgrupos resultantes
queda compuesto por una o dos observaciones, caso en el que las pruebas de
significancia basadas en aproximaciones asintóticas, como Kruskal-Wallis,
pierden validez y sus resultados deben leerse como puramente descriptivos y no
como evidencia inferencial. La significancia estadística de cada subgrupo se
evalúa, no obstante, con las mismas pruebas del Nivel 1
(Sección~\ref{sec:validacion-grupos}), con el fin de distinguir qué
subdivisiones cuentan con sustento estadístico razonable de cuáles son
artefactos de la granularidad, reportando siempre el tamaño de cada subgrupo
junto con el resultado de la prueba.

% ============================================================================
\section{Evaluación}
\label{sec:evaluacion}
% ============================================================================

Esta sección corresponde a la etapa de evaluación de CRISP-DM y abarca la
validación estadística de los grupos, la selección de variables relevantes,
el contraste con la clasificación administrativa del MINSAL y la detección de
hospitales atípicos.


\subsection{Caracterización estadística post-hoc de los grupos}
\label{sec:validacion-grupos}

Una vez obtenida la partición, los grupos se caracterizan mediante contrastes
estadísticos aplicados a variables expresadas en su escala original e
interpretable. Estos análisis se realizan entre los cuatro grupos del Nivel~1
y, de forma separada, entre los subgrupos obtenidos en el Nivel~2.

Para la caracterización se contrastan 28 variables numéricas: las 27 variables
directas utilizadas en la matriz de agrupamiento y \texttt{peso\_medio\_cma},
indicador que se conserva con fines descriptivos, pero que no participa en la
formación de los clústeres por su ausencia estructural en establecimientos sin
actividad CMA. Las ocho componentes principales de casuística utilizadas en el
clustering no se someten a este contraste, pues no poseen una interpretación
clínico-operativa directa. Adicionalmente, se examina la asociación entre la
pertenencia a un grupo y el Servicio de Salud de cada establecimiento.

La elección de la prueba depende de la naturaleza de la variable:

\begin{itemize}
  \item \textbf{Variables numéricas}: se utiliza la prueba de
    Kruskal--Wallis, que contrasta si las distribuciones de una variable difieren
    entre los grupos sin requerir normalidad. Cuando las formas de las
    distribuciones son comparables, sus resultados pueden interpretarse
    principalmente como diferencias de localización entre grupos.

\item \textbf{Variable categórica}: para la pertenencia a Servicio de Salud
  se utiliza la prueba de chi-cuadrado de independencia. Debido al número
  reducido de hospitales y a la dispersión de la tabla de contingencia, este
  resultado debe interpretarse de forma descriptiva.
\end{itemize}

El nivel de significancia se fija en $\alpha=0{,}05$. Dado que se realizan 28
contrastes numéricos en cada nivel de análisis, los $p$-valores de
Kruskal--Wallis se ajustan mediante el procedimiento de
Benjamini--Hochberg~\citep{Benjamini1995}, con el fin de controlar la tasa de
falsos descubrimientos dentro de cada familia de pruebas. La corrección no se
aplica al contraste de Servicio de Salud, pues corresponde a una única prueba
categórica por nivel.

La significancia estadística se complementa con tamaños de efecto. Para las
variables numéricas se emplea el épsilon cuadrado
($\varepsilon^2$), calculado a partir del estadístico de Kruskal--Wallis:

\begin{equation}
  \varepsilon^2 = \frac{H - (k - 1)}{n - k},
  \label{eq:epsilon2}
\end{equation}

donde $H$ corresponde al estadístico de Kruskal--Wallis, $k$ al número de
grupos y $n$ al número de hospitales analizados. Las variables se ordenan
principalmente según este tamaño de efecto, y no solo según su $p$-valor
ajustado, para jerarquizar los atributos de acuerdo con la magnitud de las
diferencias observadas entre grupos.

Para la variable categórica se reporta la V de Cramér corregida por sesgo. Esta
corrección se utiliza porque las tablas que relacionan los grupos con los
Servicios de Salud contienen numerosas categorías respecto del tamaño muestral,
situación en la cual la V de Cramér convencional puede sobreestimar la magnitud
de la asociación. Los resultados categóricos se presentan, por tanto, como
evidencia descriptiva complementaria y no como un criterio para validar la
partición.

Es necesario precisar el alcance de estos contrastes. De las 28 variables
numéricas examinadas, 27 participaron directamente en la construcción de la
matriz utilizada por Ward. Por ello, encontrar diferencias entre los grupos en
esas mismas variables no constituye una validación estadística independiente de
la partición: el algoritmo formó los grupos precisamente a partir de patrones de
similitud y diferencia en ese espacio de características. La corrección de
Benjamini--Hochberg controla los falsos descubrimientos asociados a las
comparaciones múltiples, pero no elimina la reutilización de los datos para
formar y describir los grupos.

En consecuencia, las pruebas por variable se reportan como una
caracterización descriptiva \textit{post-hoc}. Su finalidad es
identificar qué atributos presentan mayor contraste entre los perfiles ya
formados, ordenar su importancia relativa mediante tamaños de efecto y apoyar
la interpretación clínico-operativa de cada grupo. No se interpretan como
evidencia de que los clústeres constituyan la única estructura posible ni como
una validación externa de su existencia.

La evaluación de la estabilidad y del alcance de la partición se aborda mediante
evidencia complementaria: análisis de sensibilidad frente a algoritmos de
mecanismo distinto, estabilidad bajo submuestreo y una prueba de permutación del
flujo completo. Estas estrategias se describen en las
Secciones~\ref{sec:sensibilidad-analisis}, \ref{sec:bootstrap} y
\ref{sec:permutacion-circularidad}, respectivamente.


\subsubsection{Prueba de permutación sobre el flujo completo}
\label{sec:permutacion-circularidad}

Para cuantificar, en lugar de solo señalar, cuánta de la significancia
observada podría atribuirse al mero acto de agrupar --sin que exista ninguna
estructura real que separar-- se diseñó una prueba de permutación sobre el
pipeline completo. Se generan matrices institucionales nulas permutando de
forma independiente cada una de las 35 columnas entre los 65 hospitales
(\textit{shuffle} por columna), lo que destruye cualquier estructura de
covarianza real entre variables y entre hospitales, preservando exactamente
la distribución marginal (escala y forma) de cada variable individual. Sobre
cada matriz nula se repite el flujo íntegro: preprocesamiento (log1p y
escalado robusto), clustering aglomerativo Ward con $K=4$, y Kruskal-Wallis
con corrección BH sobre las mismas 28 variables interpretables. El
procedimiento se repite 500 veces, registrando en cada repetición cuántas de
las 28 variables resultan significativas ($p_{\text{BH}} \leq 0{,}05$).

Si el flujo completo --y no la estructura real de los datos-- fuera el
principal responsable de producir variables significativas, la distribución
nula debería concentrarse en valores cercanos al resultado real (26 de 28).
El resultado obtenido, reportado en el Capítulo~\ref{cap:resultados}, muestra
lo contrario: bajo la hipótesis nula de ausencia de estructura, el número de
variables significativas es sistemáticamente bajo. Esto acota, aunque no
elimina por completo, la preocupación de circularidad: demuestra que Ward no
\textit{fabrica} 24 variables significativas de la nada por el solo hecho de
particionar datos sin estructura real, aunque no demuestra que la partición
identificada sea la única o la mejor estructura posible en los datos reales,
ni sustituye a una validación externa genuina sobre una muestra o un período
distinto al empleado para construir los grupos, algo que este trabajo no
realizó.


\subsection{Validación de estabilidad mediante submuestreo}
\label{sec:bootstrap}

La robustez de la partición se evalúa mediante submuestreo repetido. En cada
una de las 100 iteraciones se selecciona aleatoriamente, sin reposición, el
80\,\% de los hospitales ($n_{\text{sub}} = 52$ de 65) y sobre ese subconjunto
se re-ejecuta el clustering aglomerativo Ward con $K=4$. El uso de submuestreo
sin reposición, en lugar del \textit{bootstrap} clásico con reposición, evita
que un mismo hospital aparezca múltiples veces en una réplica, lo que
distorsionaría las distancias euclidianas entre observaciones y produciría
estimaciones artificialmente optimistas de la cohesión. Con $n = 65$, mantener
el 80\,\% garantiza subconjuntos de tamaño suficiente para que el clustering sea
estable, mientras que el 20\,\% restante actúa como perturbación que revela si
la estructura depende de hospitales individuales influyentes.

Recalcular únicamente el coeficiente de Silhouette en cada submuestra, como se
planteó en una versión anterior de este diseño, no evalúa directamente la
estabilidad de la \emph{pertenencia} de cada hospital a su grupo: dos
submuestras podrían producir un Silhouette similar con asignaciones de
hospitales completamente distintas entre sí, sin que la sola cohesión interna
lo revele. Por ello, además del Silhouette de cada submuestra y su intervalo
de confianza empírico (percentiles 2{,}5 y 97{,}5), el procedimiento evalúa
explícitamente la estabilidad de la partición mediante cuatro métricas
adicionales:

\begin{itemize}
  \item \textbf{Concordancia global}: el Índice de Rand Ajustado (ARI) y la
    Información Mutua Normalizada (NMI) entre la partición de cada submuestra
    y la partición original de referencia, restringidos a los hospitales en
    común, cuantifican si la estructura de cuatro grupos se reproduce como un
    todo y no solo su cohesión interna.
  \item \textbf{Matriz de coasignación}: para cada par de hospitales que son
    muestreados juntos en una misma iteración, se registra si quedan en el
    mismo clúster. La proporción de iteraciones en que ambos coinciden,
    acumulada sobre las 100 réplicas, produce una matriz de $65 \times 65$
    que resume qué pares de hospitales tienden a agruparse de forma
    consistente.
  \item \textbf{Estabilidad individual}: a partir de la matriz de
    coasignación, se define la estabilidad de un hospital como su
    coasignación promedio con el resto de los integrantes de su propio
    clúster original, es decir, la proporción de sus ``vecinos'' de grupo que
    conserva a través de las submuestras.
  \item \textbf{Jaccard e estabilidad por grupo}: para cada uno de los cuatro
    clústeres originales, se calcula en cada submuestra el índice de Jaccard
    entre sus miembros presentes en esa submuestra y el clúster de la
    partición submuestreada con mayor solapamiento, y se promedia sobre las
    submuestras en que el clúster tiene representantes. Esto permite
    identificar si alguno de los cuatro grupos es sistemáticamente menos
    estable que los demás, en lugar de reportar una única cifra agregada para
    toda la partición.
\end{itemize}

Los resultados de este procedimiento, incluida la comparación entre grupos y
la identificación de los hospitales individualmente menos estables, se
presentan en la Sección~\ref{sec:res-estabilidad} del Capítulo~\ref{cap:resultados}.

\subsection{Selección de variables relevantes}
\label{sec:seleccion-variables}

La selección de características identifica qué atributos resultan determinantes
para distinguir los perfiles hospitalarios y cuáles aportan información redundante
o ruido. Para abordar este objetivo se emplean dos técnicas complementarias con
propósitos distintos: la prueba de Kruskal-Wallis, que opera directamente sobre
la partición empírica, y regresión logística con penalización LASSO, que opera
sobre la clasificación administrativa del MINSAL como referencia externa. Ambas
se aplican después de obtenida la partición para no sesgar el agrupamiento.

\subsubsection{Discriminación mediante Kruskal-Wallis}
\label{sec:kruskal-seleccion}

La prueba de Kruskal-Wallis evalúa si las distribuciones de cada variable difieren
significativamente entre los grupos obtenidos, sin asumir normalidad. Una variable
con $p \leq 0{,}05$ discrimina entre los perfiles y es relevante para
caracterizarlos, mientras que una con $p > 0{,}05$ es homogénea entre grupos y no aporta a la
diferenciación. Las variables discriminantes se ordenan por su $\varepsilon^2$
(Ecuación~\ref{eq:epsilon2}), que jerarquiza los atributos según la magnitud
real de la diferencia entre grupos, con corrección de Benjamini-Hochberg sobre
los $p$-valores, y permite construir la narrativa funcional de cada perfil.


\subsubsection{Regresión logística con penalización LASSO}
\label{sec:lasso}

La penalización $L_1$ se aplica sobre un modelo de regresión logística usando
como variable objetivo la clasificación administrativa del MINSAL. El propósito no es predecir esa etiqueta sino
identificar qué variables de actividad clínica permiten reconstruirla y, por
contraste, cuáles aportan información que la etiqueta administrativa no captura.
La penalización $L_1$ induce la dispersión de coeficientes, forzando a cero los coeficientes
de los predictores poco informativos, como se describió en la
Sección~\ref{sec:reduccion}.

Para evitar que la selección dependa de una partición particular de los datos,
el procedimiento se estabiliza mediante validación cruzada estratificada de 5
pliegues. En cada pliegue se ajusta un modelo LASSO y se registra qué variables
reciben coeficiente no nulo, la frecuencia de selección de cada variable
a lo largo de los pliegues mide su estabilidad

\begin{equation}
f_j = \frac{1}{P} \sum_{p=1}^{P}
      \mathbf{1}\!\left[\,\lvert \beta_j^{(p)} \rvert > 0 \,\right]
\label{eq:frecuencia-lasso}
\end{equation}

donde $P = 5$ es el número de pliegues, $\beta_j^{(p)}$ el coeficiente de la
variable $j$ en el pliegue $p$, y $\mathbf{1}[\cdot]$ la función indicadora. Una
variable con $f_j$ cercana a 1 es seleccionada de manera consistente y se
considera robusta, una con $f_j$ baja es inestable. El parámetro de
regularización se fija en $C = 0{,}5$ nivel moderado que equilibra la fuerza
de la penalización sin descartar variables informativas ni perder el efecto de
selección.

La combinación de ambas técnicas permite distinguir variables que discriminan
la partición empírica (Kruskal-Wallis significativo) de las que reproducen la
etiqueta administrativa (alta frecuencia LASSO), y de aquellas que aparecen en
ambas fuentes. Los resultados se presentan en el Capítulo~\ref{cap:resultados}.

\subsection{Contraste con la clasificación administrativa del MINSAL}
\label{sec:contraste-minsal}

Para cuantificar la relación entre la segmentación empírica y la clasificación
administrativa vigente, se construye la tabla de contingencia entre ambas
particiones y se calculan dos índices de concordancia: el Índice de Rand
Ajustado (ARI) y la Información Mutua Normalizada (NMI). Ambos
corrigen por el acuerdo esperado al azar y toman valor 0 cuando la concordancia
es equivalente al azar y 1 cuando las particiones son idénticas.

El propósito del contraste no es validar el clustering mediante la etiqueta
MINSAL ni demostrar que uno supera al otro, sino cuantificar en qué medida ambas
perspectivas coinciden o divergen. La clasificación MINSAL se basa en
infraestructura física y dotación de camas, mientras que la segmentación empírica se basa en
casuística clínica registrada. Son dimensiones distintas del mismo objeto, por lo
que una concordancia baja es esperable y no constituye evidencia en contra de
ninguna de las dos. Lo que un ARI bajo sí permite afirmar es que ambas
clasificaciones no son equivalentes: existen hospitales que la etiqueta oficial
trata como iguales pero que la casuística distingue, y viceversa. Lo que ARI y
NMI \textbf{no} permiten afirmar, sin embargo, es que la segmentación empírica
sea más homogénea, más precisa o mejor que MINSAL: ambos índices son medidas de
concordancia entre particiones, no de homogeneidad interna relativa. Para
evaluar directamente la hipótesis de trabajo, que sí formula una comparación
de homogeneidad, se requiere una comparación distinta a la sola concordancia.

\subsubsection{Comparación directa de homogeneidad interna}
\label{sec:comparacion-homogeneidad}

Para contrastar la hipótesis de que la segmentación funcional presenta
``mayor homogeneidad interna... que la clasificación administrativa vigente''
(Sección~\ref{sec:hipotesis}) se compara, sobre la misma matriz preprocesada
$X$ usada para el clustering y las mismas 27 variables directas, la dispersión
intragrupo de la partición Ward y la de la clasificación MINSAL, mediante tres
medidas complementarias:

\begin{itemize}
  \item \textbf{Dispersión intragrupo normalizada} (\textit{WSS/TSS}): la suma
    de cuadrados intragrupo dividida por la suma de cuadrados total de $X$.
    Un valor menor indica mayor homogeneidad interna respecto de la variación
    total de la red. A diferencia del coeficiente de Silhouette, que
    considera también la separación entre grupos, esta medida aísla
    exclusivamente la cohesión interna, que es la dimensión que la hipótesis
    contrasta.
  \item \textbf{MAD e IQR agregados por variable}: para cada una de las 27
    variables directas, se calcula la mediana de la desviación absoluta
    respecto a la mediana (MAD) y el rango intercuartílico (IQR) dentro de
    cada grupo, ponderados por su tamaño y normalizados por el MAD/IQR global
    de la variable en toda la red. Estas medidas, a diferencia de la
    desviación estándar, son robustas ante los grupos de tamaño reducido que
    integran ambas particiones.
  \item \textbf{Silhouette sobre etiquetas comparables}: el coeficiente de
    Silhouette de la clasificación MINSAL se calcula sobre el mismo espacio
    $X$ y la misma métrica euclidiana empleados para Ward, lo que permite una
    comparación directa entre ambas particiones y no solo entre una partición
    y su propia referencia interna.
\end{itemize}

Dado que MINSAL define únicamente dos categorías en el universo de 65
hospitales (Alta y Mediana Complejidad) mientras que la partición Ward
principal define cuatro, comparar estas métricas sin controlar el número de
grupos favorece mecánicamente a la solución más fragmentada: como se muestra
en la Sección~\ref{sec:nivel1}, la propia partición Ward con $K=2$ ya optimiza
las tres métricas de validación interna por encima de $K=4$. Por ello, la
comparación se reporta en dos versiones: una \textbf{no controlada}, entre la
partición Ward $K=4$ efectivamente utilizada en este trabajo y MINSAL ($K=2$),
que refleja la comparación de interés sustantivo pero con distinto número de
grupos, y una \textbf{controlada por $K$}, entre una partición Ward restringida
a $K=2$ --el mismo número de categorías que MINSAL-- y MINSAL, sobre las
mismas variables y métricas, que aísla el efecto de la homogeneidad de la
partición del efecto mecánico del número de grupos.

Finalmente, para determinar si la dispersión observada en cada partición es
menor de lo esperable por azar, se ejecuta una prueba de permutación: se
generan 2.000 reasignaciones aleatorias de los 65 hospitales que preservan
los tamaños de grupo de cada partición real (por ejemplo, 55 y 10 para
MINSAL), y se compara la dispersión normalizada observada contra esa
distribución nula. Esto permite reportar, para cada partición por separado,
si su nivel de homogeneidad interna es estadísticamente distinguible de una
agrupación aleatoria de igual tamaño de grupos, y no solo comparar los
valores puntuales entre ambas. Los resultados de esta comparación se
presentan en la Sección~\ref{sec:res-minsal} del Capítulo~\ref{cap:resultados}.

Lo que respalda la utilidad descriptiva de la solución empírica no es una
comparación aislada con MINSAL, sino el conjunto de evidencia presentado en
este capítulo: la cohesión interna medida por Silhouette, Calinski-Harabasz y
Davies-Bouldin (Sección~\ref{sec:nivel1}), la estabilidad frente a
submuestreo (Sección~\ref{sec:bootstrap}) y la caracterización posterior de
las diferencias entre grupos (Sección~\ref{sec:validacion-grupos}), con las
salvedades de circularidad ya señaladas. El contraste interno con MINSAL aporta
información descriptiva sobre la compactación de las particiones en el espacio
completo, pero no constituye una prueba confirmatoria porque Ward se ajustó en
ese mismo espacio.

\subsubsection{Validación cruzada por bloques de variables}
\label{sec:comparacion-homogeneidad-cruzada}

Para evitar esa reutilización, se añade una evaluación cruzada por bloques de
variables. La matriz se divide en las 27 variables clínico-operativas directas y
las 8 componentes PCA de casuística. En una primera dirección, Ward se ajusta
solo con las variables directas y su partición se evalúa, junto con la
clasificación MINSAL, exclusivamente sobre las componentes de casuística. En la
segunda dirección, Ward se ajusta con las componentes y ambas particiones se
evalúan sobre las variables directas. Cada bloque de evaluación se transforma
con el mismo preprocesamiento del pipeline, ajustado dentro del bloque, pero no
interviene en la construcción de las etiquetas Ward.

Se reporta WSS/TSS y Silhouette en el bloque retenido. Además, la diferencia
$\mathrm{WSS/TSS}_{\mathrm{MINSAL}}-\mathrm{WSS/TSS}_{\mathrm{Ward}}$ se
acompaña de un intervalo percentilar bootstrap del 95\,\% obtenido con 5.000
remuestreos de hospitales. El análisis se ejecuta con Ward $K=4$, que conserva
la granularidad de la solución principal, y con Ward $K=2$, para contrastar con
el mismo número de categorías de MINSAL. Una diferencia positiva cuyo intervalo
no incluye cero indica menor dispersión de Ward en un espacio que no utilizó
para agrupar. Como los bloques proceden de los mismos hospitales y las
componentes PCA se estimaron sobre la muestra completa, esta es una validación
interna por bloques y no una validación temporal o externa independiente.

\subsection{Detección de hospitales atípicos}
\label{sec:outliers-metodo}

De las técnicas discutidas en la Sección~\ref{sec:anomalias}, LOF y la distancia
de Mahalanobis fueron descartados por razones operativas. Con $n = 65$, LOF es
sensible a la elección del parámetro de vecindad y puede producir resultados
inestables en muestras pequeñas. La distancia de Mahalanobis requiere estimar la
matriz de covarianza, operación mal condicionada cuando el número de variables
(35) es del mismo orden que el número de observaciones. En su lugar se combina
una sensibilidad previa mediante Isolation Forest con dos criterios posteriores
que requieren una partición ya construida.

Antes de ajustar Ward, Isolation Forest se entrena una vez sobre los 65
hospitales y las 35 variables preprocesadas. Con contaminación $10\,\%$, sus
marcas se excluyen temporalmente y Ward $K=4$ se re-estima sobre los hospitales
restantes. La concordancia ARI/NMI con la solución base se calcula tanto sobre
los establecimientos retenidos como sobre los 65 hospitales, reasignando cada
marca al centroide más próximo de la solución entrenada sin ella. Esta exclusión
es una prueba de sensibilidad, no una regla de depuración permanente: si retirar
los casos degrada o transforma materialmente la estructura, estos se conservan
en la partición reportada y se interpretan como perfiles poco frecuentes o
influyentes, no como errores automáticos.

Después del clustering, se aplica una estrategia que combina tres criterios
complementarios, cada uno sensible a una faceta distinta de la atipicidad:

\begin{enumerate}
  \item \textbf{Coeficiente de Silhouette por observación}: un hospital con
    Silhouette negativo está más cerca del grupo vecino que del propio, señal de
    asignación ambigua.
  \item \textbf{Isolation Forest}~\citep{Liu2008}: detector basado en la facilidad
    de aislar una observación mediante particiones aleatorias del espacio de
    características. A diferencia de LOF, no requiere definir un radio de
    vecindad y es robusto con muestras pequeñas. El modelo se ajusta \textbf{una
    sola vez, de forma global}, sobre los 65 hospitales y las 35 variables de la
    matriz institucional, y no se ajusta por separado dentro de cada clúster. Se
    reporta con este alcance explícito para no sugerir que la técnica detecta
    atipicidad ``dentro de'' cada grupo cuando, en realidad, compara a cada
    hospital contra la totalidad de la red.
  \item \textbf{Distancia euclidiana al centroide}: se calcula la distancia de
    cada hospital al centroide de su propio clúster y se normaliza mediante la
    mediana de la desviación absoluta (MAD) de esa distancia, calculada sobre el
    conjunto completo de 65 hospitales y no dentro de cada clúster por separado.
    Se marca como atípico todo hospital cuya distancia normalizada supere un
    valor de referencia equivalente a 3 desviaciones estándar bajo normalidad
    ($\text{MAD} \times 1{,}4826$). Esta normalización global, en lugar de una
    desviación estándar estimada dentro de cada grupo, evita depender de una
    dispersión intra-clúster calculada con apenas 2 o 3 observaciones, como
    ocurre en $C2$ ($n=2$) y $C0$ ($n=3$), donde esa estimación sería
    estadísticamente inestable.
\end{enumerate}

El parámetro de contaminación de Isolation Forest ($10\,\%$) no se deriva de
una regla formal ni de un umbral estadístico: es un valor de uso frecuente en
la literatura aplicada de detección de anomalías, adoptado aquí como punto de
partida operativo, y su elección determina de forma directa cuántas
observaciones el algoritmo está obligado a señalar como atípicas,
independientemente de la estructura real de los datos. Por esta razón se
evaluó explícitamente la sensibilidad del resultado a este parámetro,
repitiendo la detección con niveles de contaminación de $5\,\%$, $10\,\%$,
$15\,\%$ y $20\,\%$, y reportando qué hospitales son detectados de forma
consistente en todos los niveles explorados, en lugar de presentar el
resultado de un único valor como si fuera una propiedad emergente de los
datos. Los resultados de esta sensibilidad se presentan en el
Capítulo~\ref{cap:resultados}.

Un hospital se clasifica como atípico si lo señala al menos uno de los
tres criterios, y como \textbf{atípico robusto} si lo señalan dos o más. Esta
distinción es metodológicamente importante, ya que los tres criterios capturan fenómenos
distintos (asignación ambigua, baja densidad global, distancia extrema al
propio centroide), por lo que su intersección identifica los casos con mayor
consistencia entre criterios, mientras que los
marcados por un solo criterio corresponden a hospitales en la frontera entre
perfiles. La condición de ``atípico'' es, en todos los casos, puramente
estadística y descriptiva: no implica por sí sola una anomalía en el sentido
de un error, una falla de gestión o una necesidad de auditoría, y su
interpretación sustantiva --especialización genuina, perfil híbrido o posible
problema de codificación-- se desarrolla caso por caso en el
Capítulo~\ref{cap:resultados}.

\subsection{Interpretabilidad y proyecciones visuales}
\label{sec:interpretabilidad-metodo}

La descripción clínica de cada grupo se apoya en una interpretación
\textit{post-hoc} basada en las medianas de las variables más discriminantes,
realizada exclusivamente después del clustering para no sesgar la formación de
los grupos. Complementariamente, se proyecta la matriz a dos dimensiones
mediante cinco técnicas independientes: PCA, UMAP, MDS, t-SNE e Isomap. El
propósito de estas proyecciones es exclusivamente ilustrativo y exploratorio:
permiten inspeccionar de forma visual si la separación entre grupos observada
bajo PCA se mantiene reconocible bajo otras técnicas de reducción con supuestos
geométricos distintos, lo que descarta que esa separación sea un artefacto
exclusivo de la proyección PCA en particular. La coincidencia visual entre
proyecciones no constituye, sin embargo, evidencia cuantitativa de la validez
del clustering ni sustituye a los índices de concordancia de partición (ARI,
NMI) que se emplean para ese propósito en la Sección~\ref{sec:sensibilidad}: una
proyección bidimensional descarta necesariamente información de las 35
variables originales, y la similitud visual entre dos nubes de puntos es una
apreciación cualitativa sujeta a la interpretación de quien observa la figura,
no una medida estadística.

% ============================================================================
\section{Análisis de sensibilidad del diseño}
\label{sec:sensibilidad-diseno}
% ============================================================================

Las secciones anteriores documentan cuatro decisiones de diseño que no se derivan
de una regla única e indiscutible: el algoritmo de agrupamiento
(Sección~\ref{sec:algoritmo}), el criterio de elegibilidad de hospitales
(Sección~\ref{sec:elegibilidad}), la agregación de los seis años en un único
conjunto de egresos (Sección~\ref{sec:variables}) y el tratamiento de los
vectores de casuística como datos euclidianos en lugar de composicionales
(Sección~\ref{sec:pca}). Esta sección define el procedimiento con que se evalúa
cuánto dependen los resultados de cada una de ellas. Los resultados de las cuatro
pruebas se reportan en el Capítulo~\ref{cap:resultados}.

Las tres últimas pruebas comparten un mismo protocolo. Se altera una sola decisión
de diseño, se recalcula el pipeline completo --construcción de variables,
reducción PCA, preprocesamiento y clustering aglomerativo Ward con $K=4$-- y se
compara la partición obtenida con la del pipeline definitivo mediante el índice de
Rand ajustado (ARI) y la información mutua normalizada (NMI), acompañados del
coeficiente de Silhouette como medida de cohesión interna de cada solución. La
elección de ARI y NMI responde a que ambos índices son invariantes a la
permutación de etiquetas: la numeración que produce el algoritmo carece de
contenido semántico y depende del orden en que se construyen las fusiones del
árbol, de modo que dos particiones equivalentes pueden designar el mismo grupo con
etiquetas distintas. Por esta razón, la comparación entre particiones se realiza
con índices invariantes y no contabilizando coincidencias de etiquetas, y la
interpretación de cada prueba distingue explícitamente entre renumeración de
clústeres y reasignación sustantiva de establecimientos.

Las cuatro pruebas constituyen un análisis de sensibilidad interna: las
particiones alternativas se construyen sobre la misma fuente de datos y el mismo
conjunto de variables que la solución definitiva, de modo que evalúan la
dependencia del resultado respecto de decisiones metodológicas y no su validez
frente a una fuente externa. La única prueba que altera el universo de
establecimientos es la del criterio de elegibilidad, y en ese caso la comparación
se restringe a los 65 hospitales presentes en ambas soluciones.

\subsection{Análisis de sensibilidad algorítmica}
\label{sec:sensibilidad}

Para evaluar si la estructura identificada es una propiedad genuina de los datos
y no un artefacto del algoritmo aglomerativo Ward, se contrasta el resultado con
cinco familias adicionales de técnicas de clustering, resumidas en la
Tabla~\ref{tab:sensibilidad}.

\begin{table}[htbp]
\centering
\caption{Técnicas de clustering empleadas en el análisis de sensibilidad.}
\label{tab:sensibilidad}
\begin{tabular}{p{4.2cm} p{3.0cm} p{4.5cm}}
\toprule
\textbf{Técnica} & \textbf{Parámetros explorados} & \textbf{Propósito} \\
\midrule
K-means++ & $k \in \{2, \ldots, 8\}$ y $\{5,\ldots,20\}$ & Contraste particional. \\
\addlinespace
Mezclas gaussianas (GMM) & $k \in \{2, \ldots, 8\}$, BIC & Contraste probabilístico. \\
\addlinespace
HDBSCAN & \texttt{min\_cluster\_size} $\in \{3,\ldots,8\}$ & Descubre $k$ por densidad. \\
\addlinespace
MST-kNN & $k$ de vecindad $=6$ & Algoritmo del depto.\ USACH,
descubre $k$ y observaciones aisladas. \\
\addlinespace
Proyecciones 2D & PCA, UMAP, MDS, t-SNE, Isomap & Inspección visual
exploratoria de la separación, no constituye validación cuantitativa. \\
\bottomrule
\end{tabular}
\end{table}

La inclusión del algoritmo \textbf{MST-kNN}~\citep{inostroza2008}, utilizado previamente en el ámbito hospitalario nacional por \citet{Villalobos2016}, es relevante por dos
razones: es el algoritmo de referencia desarrollado en el Departamento de
Ingeniería Informática de la USACH, lo que sitúa este trabajo en continuidad con
la línea de investigación local, y no requiere especificar el número de grupos,
descubriendo automáticamente tanto los clusters como los hospitales de perfil
único. Adicionalmente, se exploran
particiones con $k \in \{5, \ldots, 20\}$ para verificar el comportamiento de las
métricas en rangos de alta granularidad y una gran cantidad de grupos.

\subsection{Sensibilidad al criterio de elegibilidad}
\label{sec:sensibilidad-elegibilidad}

La segunda prueba evalúa el efecto del criterio de elegibilidad sobre la
partición final y no solo sobre el error de los indicadores individuales, que ya
se cuantificó en la Sección~\ref{sec:elegibilidad}. Como se estableció allí, no
existen hospitales con 3, 4 o 5 años de datos, por lo que elevar el umbral a 4 o
5 años no altera el universo analizado y no constituye una prueba informativa. La
perturbación que sí modifica el universo es la contraria: eliminar por completo
el filtro de presencia temporal e incorporar los 7 hospitales excluidos
(Tabla~\ref{tab:descartados}) con los datos parciales de uno o dos años que
efectivamente tienen disponibles. Se reconstruye la matriz institucional sobre
este universo ampliado de 72 hospitales, se repite el clustering aglomerativo
Ward con $K=4$ y la concordancia con la partición original se calcula
restringiendo la comparación a los 65 hospitales presentes en ambas soluciones.

\subsection{Sensibilidad a la exclusión de los años de pandemia}
\label{sec:sensibilidad-temporal}

La tercera prueba evalúa si la agregación de los seis años en un único conjunto
de egresos, sin ponderación temporal (Sección~\ref{sec:variables}), oculta
cambios estructurales asociados a la pandemia de COVID-19. Para ello se recalcula
la matriz institucional completa excluyendo los años 2020 y 2021, manteniendo
fijo el universo de 65 hospitales elegibles definido en la
Sección~\ref{sec:elegibilidad}. La elegibilidad no se re-deriva sobre el período
reducido, de modo que cualquier diferencia observada sea atribuible al cambio de
período y no a un cambio en la composición del universo de análisis. Sobre esta
matriz alternativa, construida con 4{,}158{,}029 egresos de los años 2019 y
2022--2024 --un 71{,}6\,\% del total del período--, se repiten el
preprocesamiento de la Sección~\ref{sec:preprocesamiento} y el clustering
aglomerativo Ward con $K=4$.

Esta prueba incorpora un segundo nivel de análisis, además de la comparación
entre particiones. Para cada variable de la matriz institucional se calcula el
cambio relativo mediano entre hospitales al pasar del período completo al período
reducido, junto con la correlación entre ambas versiones de la variable. El
propósito es distinguir si una eventual estabilidad de la partición se sostiene
sobre indicadores igualmente estables o si, por el contrario, coexiste con
variables individualmente sensibles al cambio de período.

\subsection{Sensibilidad a la transformación composicional}
\label{sec:sensibilidad-composicional}

La cuarta prueba evalúa si el problema de clausura composicional descrito en la
Sección~\ref{sec:pca} afecta materialmente los resultados. Para ello se recalcula
la representación de casuística aplicando una transformación CLR
(\textit{centered log-ratio}) a los tres vectores de proporciones antes del PCA,
en lugar de operar sobre las proporciones brutas. La transformación CLR proyecta
cada composición al espacio logarítmico de razones respecto a su media
geométrica, lo que remueve la restricción de suma constante y constituye la
alternativa estándar en el análisis de datos composicionales. Los ceros
estructurales de los vectores de proporciones --frecuentes en establecimientos
que no registran actividad en determinado capítulo, sección o GRD-- se tratan
mediante un pseudo-conteo multiplicativo ($\delta = 1/(2D)$, con $D$ el número de
partes de cada composición), que evita la indeterminación de $\log(0)$ sin
distorsionar de forma relevante las proporciones no nulas.

Sobre la representación CLR se aplica PCA con el mismo número de componentes
($K=8$) y se reconstruye la matriz institucional completa, sustituyendo
únicamente las componentes de casuística y dejando intactas las 27 variables
directas. Como verificación adicional, la selección de la variante de capítulos
(principal frente a ponderada) realizada en la Sección~\ref{sec:pca-detalle} con
el criterio no circular de varianza explicada se repite sobre la representación
CLR, para comprobar que esa decisión de diseño no depende de la geometría
empleada.

Cabe anticipar por qué la geometría euclidiana sobre las componentes PCA de las
proporciones brutas puede resultar razonable aun sin ser la elección teóricamente
canónica para datos composicionales. El PCA reduce las 63 proporciones originales
a 8 componentes antes de calcular cualquier distancia, y esa reducción atenúa,
aunque no elimina, la distorsión de clausura, porque las componentes resultantes
capturan patrones de covariación conjunta entre partes en lugar de operar sobre
cada proporción aislada. Esto no descarta el problema: el vector de Top-20 GRD,
en particular, incluye proporciones individuales de hasta $0{,}978$ en algunos
hospitales altamente concentrados en un único GRD --precisamente el tipo de
composición donde el efecto de clausura es más severo--, por lo que la robustez
frente a este problema no puede darse por sentada a priori y requiere la
comparación empírica que esta prueba provee. El pipeline no explora la
transformación ILR (\textit{isometric log-ratio}), que evita además la
singularidad de la matriz de covarianza introducida por CLR, y que constituye una
extensión natural de este análisis para trabajo futuro.

% ============================================================================
\section{Despliegue}
\label{sec:despliegue}
% ============================================================================

La etapa de despliegue en CRISP-DM corresponde a la forma en que los resultados
del análisis se comunican y quedan disponibles para su uso. En el contexto de
esta investigación, el despliegue no implica la integración de un sistema
automático en producción sino la presentación de los perfiles funcionales en un
formato interpretable y utilizable por gestores sanitarios.

El producto principal del despliegue son los perfiles funcionales de los
hospitales públicos chilenos de alta y mediana complejidad, presentados en el
Capítulo~\ref{cap:resultados}. Cada perfil se describe mediante las medianas de
las variables más discriminantes, una etiqueta funcional derivada del análisis
\textit{post-hoc} y el listado de los establecimientos que lo integran. Esta
descripción está diseñada para ser legible sin conocimiento previo del pipeline,
de modo que pueda ser adoptada por equipos de gestión sin formación en
machine learning.

El despliegue incluye además tres productos complementarios. El primero es el
contraste entre la segmentación empírica y la clasificación administrativa del
MINSAL, que identifica los establecimientos que cambian de categoría entre ambas
perspectivas y cuantifica la magnitud de la discrepancia mediante ARI y NMI. El
segundo es el listado de hospitales atípicos con su grado de singularidad
(atípico o atípico robusto) y el criterio o criterios que los señalan, información
potencialmente útil como insumo para revisión individual o auditoría. El
tercero es el pipeline completo documentado y reproducible, que permite actualizar
los perfiles cuando se incorporen nuevos años de datos GRD al sistema, cumpliendo
con el requisito de reproducibilidad establecido en la
Sección~\ref{sec:comprension-negocio}.
